% Options for packages loaded elsewhere
\PassOptionsToPackage{unicode}{hyperref}
\PassOptionsToPackage{hyphens}{url}
\PassOptionsToPackage{dvipsnames,svgnames,x11names}{xcolor}
%
\documentclass[
  letterpaper,
  DIV=11,
  numbers=noendperiod]{scrreprt}

\usepackage{amsmath,amssymb}
\usepackage{iftex}
\ifPDFTeX
  \usepackage[T1]{fontenc}
  \usepackage[utf8]{inputenc}
  \usepackage{textcomp} % provide euro and other symbols
\else % if luatex or xetex
  \usepackage{unicode-math}
  \defaultfontfeatures{Scale=MatchLowercase}
  \defaultfontfeatures[\rmfamily]{Ligatures=TeX,Scale=1}
\fi
\usepackage{lmodern}
\ifPDFTeX\else  
    % xetex/luatex font selection
\fi
% Use upquote if available, for straight quotes in verbatim environments
\IfFileExists{upquote.sty}{\usepackage{upquote}}{}
\IfFileExists{microtype.sty}{% use microtype if available
  \usepackage[]{microtype}
  \UseMicrotypeSet[protrusion]{basicmath} % disable protrusion for tt fonts
}{}
\makeatletter
\@ifundefined{KOMAClassName}{% if non-KOMA class
  \IfFileExists{parskip.sty}{%
    \usepackage{parskip}
  }{% else
    \setlength{\parindent}{0pt}
    \setlength{\parskip}{6pt plus 2pt minus 1pt}}
}{% if KOMA class
  \KOMAoptions{parskip=half}}
\makeatother
\usepackage{xcolor}
\setlength{\emergencystretch}{3em} % prevent overfull lines
\setcounter{secnumdepth}{5}
% Make \paragraph and \subparagraph free-standing
\ifx\paragraph\undefined\else
  \let\oldparagraph\paragraph
  \renewcommand{\paragraph}[1]{\oldparagraph{#1}\mbox{}}
\fi
\ifx\subparagraph\undefined\else
  \let\oldsubparagraph\subparagraph
  \renewcommand{\subparagraph}[1]{\oldsubparagraph{#1}\mbox{}}
\fi

\usepackage{color}
\usepackage{fancyvrb}
\newcommand{\VerbBar}{|}
\newcommand{\VERB}{\Verb[commandchars=\\\{\}]}
\DefineVerbatimEnvironment{Highlighting}{Verbatim}{commandchars=\\\{\}}
% Add ',fontsize=\small' for more characters per line
\usepackage{framed}
\definecolor{shadecolor}{RGB}{241,243,245}
\newenvironment{Shaded}{\begin{snugshade}}{\end{snugshade}}
\newcommand{\AlertTok}[1]{\textcolor[rgb]{0.68,0.00,0.00}{#1}}
\newcommand{\AnnotationTok}[1]{\textcolor[rgb]{0.37,0.37,0.37}{#1}}
\newcommand{\AttributeTok}[1]{\textcolor[rgb]{0.40,0.45,0.13}{#1}}
\newcommand{\BaseNTok}[1]{\textcolor[rgb]{0.68,0.00,0.00}{#1}}
\newcommand{\BuiltInTok}[1]{\textcolor[rgb]{0.00,0.23,0.31}{#1}}
\newcommand{\CharTok}[1]{\textcolor[rgb]{0.13,0.47,0.30}{#1}}
\newcommand{\CommentTok}[1]{\textcolor[rgb]{0.37,0.37,0.37}{#1}}
\newcommand{\CommentVarTok}[1]{\textcolor[rgb]{0.37,0.37,0.37}{\textit{#1}}}
\newcommand{\ConstantTok}[1]{\textcolor[rgb]{0.56,0.35,0.01}{#1}}
\newcommand{\ControlFlowTok}[1]{\textcolor[rgb]{0.00,0.23,0.31}{#1}}
\newcommand{\DataTypeTok}[1]{\textcolor[rgb]{0.68,0.00,0.00}{#1}}
\newcommand{\DecValTok}[1]{\textcolor[rgb]{0.68,0.00,0.00}{#1}}
\newcommand{\DocumentationTok}[1]{\textcolor[rgb]{0.37,0.37,0.37}{\textit{#1}}}
\newcommand{\ErrorTok}[1]{\textcolor[rgb]{0.68,0.00,0.00}{#1}}
\newcommand{\ExtensionTok}[1]{\textcolor[rgb]{0.00,0.23,0.31}{#1}}
\newcommand{\FloatTok}[1]{\textcolor[rgb]{0.68,0.00,0.00}{#1}}
\newcommand{\FunctionTok}[1]{\textcolor[rgb]{0.28,0.35,0.67}{#1}}
\newcommand{\ImportTok}[1]{\textcolor[rgb]{0.00,0.46,0.62}{#1}}
\newcommand{\InformationTok}[1]{\textcolor[rgb]{0.37,0.37,0.37}{#1}}
\newcommand{\KeywordTok}[1]{\textcolor[rgb]{0.00,0.23,0.31}{#1}}
\newcommand{\NormalTok}[1]{\textcolor[rgb]{0.00,0.23,0.31}{#1}}
\newcommand{\OperatorTok}[1]{\textcolor[rgb]{0.37,0.37,0.37}{#1}}
\newcommand{\OtherTok}[1]{\textcolor[rgb]{0.00,0.23,0.31}{#1}}
\newcommand{\PreprocessorTok}[1]{\textcolor[rgb]{0.68,0.00,0.00}{#1}}
\newcommand{\RegionMarkerTok}[1]{\textcolor[rgb]{0.00,0.23,0.31}{#1}}
\newcommand{\SpecialCharTok}[1]{\textcolor[rgb]{0.37,0.37,0.37}{#1}}
\newcommand{\SpecialStringTok}[1]{\textcolor[rgb]{0.13,0.47,0.30}{#1}}
\newcommand{\StringTok}[1]{\textcolor[rgb]{0.13,0.47,0.30}{#1}}
\newcommand{\VariableTok}[1]{\textcolor[rgb]{0.07,0.07,0.07}{#1}}
\newcommand{\VerbatimStringTok}[1]{\textcolor[rgb]{0.13,0.47,0.30}{#1}}
\newcommand{\WarningTok}[1]{\textcolor[rgb]{0.37,0.37,0.37}{\textit{#1}}}

\providecommand{\tightlist}{%
  \setlength{\itemsep}{0pt}\setlength{\parskip}{0pt}}\usepackage{longtable,booktabs,array}
\usepackage{calc} % for calculating minipage widths
% Correct order of tables after \paragraph or \subparagraph
\usepackage{etoolbox}
\makeatletter
\patchcmd\longtable{\par}{\if@noskipsec\mbox{}\fi\par}{}{}
\makeatother
% Allow footnotes in longtable head/foot
\IfFileExists{footnotehyper.sty}{\usepackage{footnotehyper}}{\usepackage{footnote}}
\makesavenoteenv{longtable}
\usepackage{graphicx}
\makeatletter
\def\maxwidth{\ifdim\Gin@nat@width>\linewidth\linewidth\else\Gin@nat@width\fi}
\def\maxheight{\ifdim\Gin@nat@height>\textheight\textheight\else\Gin@nat@height\fi}
\makeatother
% Scale images if necessary, so that they will not overflow the page
% margins by default, and it is still possible to overwrite the defaults
% using explicit options in \includegraphics[width, height, ...]{}
\setkeys{Gin}{width=\maxwidth,height=\maxheight,keepaspectratio}
% Set default figure placement to htbp
\makeatletter
\def\fps@figure{htbp}
\makeatother

\KOMAoption{captions}{tableheading}
\makeatletter
\@ifpackageloaded{tcolorbox}{}{\usepackage[skins,breakable]{tcolorbox}}
\@ifpackageloaded{fontawesome5}{}{\usepackage{fontawesome5}}
\definecolor{quarto-callout-color}{HTML}{909090}
\definecolor{quarto-callout-note-color}{HTML}{0758E5}
\definecolor{quarto-callout-important-color}{HTML}{CC1914}
\definecolor{quarto-callout-warning-color}{HTML}{EB9113}
\definecolor{quarto-callout-tip-color}{HTML}{00A047}
\definecolor{quarto-callout-caution-color}{HTML}{FC5300}
\definecolor{quarto-callout-color-frame}{HTML}{acacac}
\definecolor{quarto-callout-note-color-frame}{HTML}{4582ec}
\definecolor{quarto-callout-important-color-frame}{HTML}{d9534f}
\definecolor{quarto-callout-warning-color-frame}{HTML}{f0ad4e}
\definecolor{quarto-callout-tip-color-frame}{HTML}{02b875}
\definecolor{quarto-callout-caution-color-frame}{HTML}{fd7e14}
\makeatother
\makeatletter
\makeatother
\makeatletter
\@ifpackageloaded{bookmark}{}{\usepackage{bookmark}}
\makeatother
\makeatletter
\@ifpackageloaded{caption}{}{\usepackage{caption}}
\AtBeginDocument{%
\ifdefined\contentsname
  \renewcommand*\contentsname{Table of contents}
\else
  \newcommand\contentsname{Table of contents}
\fi
\ifdefined\listfigurename
  \renewcommand*\listfigurename{List of Figures}
\else
  \newcommand\listfigurename{List of Figures}
\fi
\ifdefined\listtablename
  \renewcommand*\listtablename{List of Tables}
\else
  \newcommand\listtablename{List of Tables}
\fi
\ifdefined\figurename
  \renewcommand*\figurename{Figure}
\else
  \newcommand\figurename{Figure}
\fi
\ifdefined\tablename
  \renewcommand*\tablename{Table}
\else
  \newcommand\tablename{Table}
\fi
}
\@ifpackageloaded{float}{}{\usepackage{float}}
\floatstyle{ruled}
\@ifundefined{c@chapter}{\newfloat{codelisting}{h}{lop}}{\newfloat{codelisting}{h}{lop}[chapter]}
\floatname{codelisting}{Listing}
\newcommand*\listoflistings{\listof{codelisting}{List of Listings}}
\makeatother
\makeatletter
\@ifpackageloaded{caption}{}{\usepackage{caption}}
\@ifpackageloaded{subcaption}{}{\usepackage{subcaption}}
\makeatother
\makeatletter
\@ifpackageloaded{tcolorbox}{}{\usepackage[skins,breakable]{tcolorbox}}
\makeatother
\makeatletter
\@ifundefined{shadecolor}{\definecolor{shadecolor}{rgb}{.97, .97, .97}}
\makeatother
\makeatletter
\makeatother
\makeatletter
\makeatother
\ifLuaTeX
  \usepackage{selnolig}  % disable illegal ligatures
\fi
\IfFileExists{bookmark.sty}{\usepackage{bookmark}}{\usepackage{hyperref}}
\IfFileExists{xurl.sty}{\usepackage{xurl}}{} % add URL line breaks if available
\urlstyle{same} % disable monospaced font for URLs
\hypersetup{
  pdftitle={Snowflake Feature Store Implementation Guide},
  pdfauthor={Simon Field},
  pdfkeywords={snowflake, feature store, ml, machine learning, mlops},
  colorlinks=true,
  linkcolor={blue},
  filecolor={Maroon},
  citecolor={Blue},
  urlcolor={Blue},
  pdfcreator={LaTeX via pandoc}}

\title{Snowflake Feature Store Implementation Guide}
\usepackage{etoolbox}
\makeatletter
\providecommand{\subtitle}[1]{% add subtitle to \maketitle
  \apptocmd{\@title}{\par {\large #1 \par}}{}{}
}
\makeatother
\subtitle{Best Practices for Building Production ML Feature Platforms}
\author{Simon Field}
\date{2026-01-21}

\begin{document}
\maketitle
\ifdefined\Shaded\renewenvironment{Shaded}{\begin{tcolorbox}[boxrule=0pt, breakable, enhanced, frame hidden, sharp corners, interior hidden, borderline west={3pt}{0pt}{shadecolor}]}{\end{tcolorbox}}\fi

\renewcommand*\contentsname{Table of contents}
{
\hypersetup{linkcolor=}
\setcounter{tocdepth}{2}
\tableofcontents
}
\bookmarksetup{startatroot}

\hypertarget{snowflake-feature-store-implementation-guide}{%
\chapter{Snowflake Feature Store Implementation
Guide}\label{snowflake-feature-store-implementation-guide}}

Best Practices for Building Production ML Feature Platforms

\hfill\break

\hypertarget{welcome}{%
\section*{Welcome}\label{welcome}}
\addcontentsline{toc}{section}{Welcome}

\markright{Welcome}

This guide provides comprehensive, practical guidance for implementing a
Feature Store using Snowflake's native Feature Store capabilities.
Whether you're building your first ML feature platform or migrating from
another solution, this guide will help you design, build, and operate a
production-grade Feature Store.

\hypertarget{what-youll-learn}{%
\section*{What You'll Learn}\label{what-youll-learn}}
\addcontentsline{toc}{section}{What You'll Learn}

\markright{What You'll Learn}

\begin{tcolorbox}[enhanced jigsaw, toprule=.15mm, colframe=quarto-callout-note-color-frame, leftrule=.75mm, rightrule=.15mm, left=2mm, breakable, colback=white, arc=.35mm, bottomrule=.15mm, opacityback=0]
\begin{minipage}[t]{5.5mm}
\textcolor{quarto-callout-note-color}{\faInfo}
\end{minipage}%
\begin{minipage}[t]{\textwidth - 5.5mm}

\textbf{Key Topics Covered}\vspace{2mm}

\begin{itemize}
\tightlist
\item
  \textbf{Core Concepts} - Entities, FeatureViews, and the Feature Store
  architecture
\item
  \textbf{Design Patterns} - Organization, naming conventions, and
  governance
\item
  \textbf{Feature Engineering} - Temporal features, aggregations, and
  pipelines
\item
  \textbf{Production Operations} - Online serving, preprocessing, and
  monitoring
\item
  \textbf{Advanced Patterns} - Streaming, CI/CD, and cross-domain ML
\end{itemize}

\end{minipage}%
\end{tcolorbox}

\hypertarget{prerequisites}{%
\section*{Prerequisites}\label{prerequisites}}
\addcontentsline{toc}{section}{Prerequisites}

\markright{Prerequisites}

Before starting, ensure you have:

\begin{itemize}
\tightlist
\item
  ✅ Snowflake account (Standard Edition or higher)
\item
  ✅ Python 3.8+ with \texttt{snowflake-ml-python} \textgreater= 1.7.0
\item
  ✅ Basic familiarity with SQL and Python
\item
  ✅ Understanding of ML concepts (training, inference, features)
\end{itemize}

\hypertarget{quick-start}{%
\section*{Quick Start}\label{quick-start}}
\addcontentsline{toc}{section}{Quick Start}

\markright{Quick Start}

\begin{Shaded}
\begin{Highlighting}[]
\CommentTok{\# Install the required package}
\CommentTok{\# pip install snowflake{-}ml{-}python}

\ImportTok{from}\NormalTok{ snowflake.snowpark }\ImportTok{import}\NormalTok{ Session}
\ImportTok{from}\NormalTok{ snowflake.ml.feature\_store }\ImportTok{import}\NormalTok{ FeatureStore, Entity, FeatureView}

\CommentTok{\# Connect to Snowflake}
\NormalTok{session }\OperatorTok{=}\NormalTok{ Session.builder.configs(connection\_params).create()}

\CommentTok{\# Initialize Feature Store}
\NormalTok{fs }\OperatorTok{=}\NormalTok{ FeatureStore(}
\NormalTok{    session}\OperatorTok{=}\NormalTok{session,}
\NormalTok{    database}\OperatorTok{=}\StringTok{"ML\_DB"}\NormalTok{,}
\NormalTok{    name}\OperatorTok{=}\StringTok{"FEATURE\_STORE"}\NormalTok{,}
\NormalTok{    default\_warehouse}\OperatorTok{=}\StringTok{"COMPUTE\_WH"}\NormalTok{,}
\NormalTok{)}

\CommentTok{\# You\textquotesingle{}re ready to start building features!}
\end{Highlighting}
\end{Shaded}

\hypertarget{how-to-use-this-guide}{%
\section*{How to Use This Guide}\label{how-to-use-this-guide}}
\addcontentsline{toc}{section}{How to Use This Guide}

\markright{How to Use This Guide}

\begin{longtable}[]{@{}
  >{\raggedright\arraybackslash}p{(\columnwidth - 2\tabcolsep) * \real{0.4783}}
  >{\raggedright\arraybackslash}p{(\columnwidth - 2\tabcolsep) * \real{0.5217}}@{}}
\toprule\noalign{}
\begin{minipage}[b]{\linewidth}\raggedright
Your Goal
\end{minipage} & \begin{minipage}[b]{\linewidth}\raggedright
Start Here
\end{minipage} \\
\midrule\noalign{}
\endhead
\bottomrule\noalign{}
\endlastfoot
New to Feature Store & \protect\hyperlink{introduction}{Chapter 0:
Introduction} \\
Design a Feature Store &
\protect\hyperlink{chapter-02-design-organization}{Chapter 2: Design \&
Organization} \\
Build temporal features &
\protect\hyperlink{chapter-06-temporal-features}{Chapter 6: Temporal
Features} \\
Enable online serving &
\protect\hyperlink{chapter-08-online-features}{Chapter 8: Online
Features} \\
Migrate from Tecton/Feast &
\protect\hyperlink{chapter-13-migration-guide}{Chapter 13: Migration
Guide} \\
\end{longtable}

\hypertarget{code-examples}{%
\section*{Code Examples}\label{code-examples}}
\addcontentsline{toc}{section}{Code Examples}

\markright{Code Examples}

All code examples in this guide are:

\begin{itemize}
\tightlist
\item
  ✅ \textbf{Tested} - Every example is verified to work
\item
  ✅ \textbf{Copy-ready} - Click the copy button on any code block
\item
  ✅ \textbf{Self-contained} - Each example includes necessary imports
\item
  ✅ \textbf{Available in repo} - Browse the \texttt{\_code/}
  directories
\end{itemize}

\hypertarget{feedback-contributions}{%
\section*{Feedback \& Contributions}\label{feedback-contributions}}
\addcontentsline{toc}{section}{Feedback \& Contributions}

\markright{Feedback \& Contributions}

Found an issue or want to suggest an improvement?

\begin{itemize}
\tightlist
\item
  📝 Use the ``Edit this page'' link on any page
\item
  🐛
  \href{https://github.com/snowcat-cloud/FeatureStore_BestPractice/issues}{Report
  an issue}
\item
  💬 Questions? Reach out via GitHub Discussions
\end{itemize}

\begin{center}\rule{0.5\linewidth}{0.5pt}\end{center}

\emph{Last updated: January 2026}

\part{Getting Started}

\hypertarget{introduction}{%
\chapter{Introduction}\label{introduction}}

Setup, prerequisites, and guide overview

\hfill\break

\hypertarget{overview}{%
\section{Overview}\label{overview}}

This guide provides comprehensive best practices for implementing
Snowflake Feature Store in production ML systems. It covers the full
lifecycle from initial design through deployment and ongoing operations.

The Snowflake Feature Store is a managed data management layer that
helps ML teams organize, store, and serve features at scale. It
provides:

\begin{itemize}
\tightlist
\item
  \textbf{Centralized Feature Management}: A single source of truth for
  feature definitions and metadata
\item
  \textbf{Automated Feature Engineering Pipelines}: Dynamic Table-based
  materialization (pre-computed) or View-based (query-time computation)
\item
  \textbf{Point-in-Time Feature Retrieval}: Temporally correct feature
  joins for ML training and batch inference, preventing data leakage
\item
  \textbf{Online Serving}: Low-latency feature retrieval via Online
  Feature Tables for real-time inference
\item
  \textbf{Governance \& Lineage}: Full tracking from source data to
  model predictions
\end{itemize}

\hypertarget{whats-covered-in-this-guide}{%
\section{What's Covered in This
Guide}\label{whats-covered-in-this-guide}}

\begin{longtable}[]{@{}
  >{\raggedright\arraybackslash}p{(\columnwidth - 4\tabcolsep) * \real{0.3103}}
  >{\raggedright\arraybackslash}p{(\columnwidth - 4\tabcolsep) * \real{0.2414}}
  >{\raggedright\arraybackslash}p{(\columnwidth - 4\tabcolsep) * \real{0.4483}}@{}}
\toprule\noalign{}
\begin{minipage}[b]{\linewidth}\raggedright
Chapter
\end{minipage} & \begin{minipage}[b]{\linewidth}\raggedright
Topic
\end{minipage} & \begin{minipage}[b]{\linewidth}\raggedright
Description
\end{minipage} \\
\midrule\noalign{}
\endhead
\bottomrule\noalign{}
\endlastfoot
\textbf{00} & Introduction & Setup, prerequisites, guide overview \\
\protect\hyperlink{core-concepts}{\textbf{01}} & Concepts & Core
objects, terminology, transformation taxonomy \\
\protect\hyperlink{chapter-02-design-organization}{\textbf{02}} & Design
\& Organization & Feature Store structure, multi-environment patterns \\
\protect\hyperlink{chapter-03-entities-hierarchies}{\textbf{03}} &
Entities \& Hierarchies & Entity keys (simple \& compound),
relationships \\
\protect\hyperlink{chapter-04-feature-views}{\textbf{04}} & Feature
Views & DT vs View, versioning, lifecycle management \\
\protect\hyperlink{chapter-05-feature-pipelines}{\textbf{05}} & Feature
Pipelines & DBT, Dynamic Tables, Temporal API pipelines \\
\protect\hyperlink{chapter-06-temporal-features}{\textbf{06}} & Temporal
Features & Point-in-time correctness, tiling, aggregations \\
\protect\hyperlink{chapter-07-aggregations-api}{\textbf{07}} &
Aggregations API & \texttt{Feature} class, windows, rollups \\
\protect\hyperlink{chapter-08-online-features}{\textbf{08}} & Online
Features & Online Feature Tables, serving patterns \\
\protect\hyperlink{chapter-09-preprocessing}{\textbf{09}} &
Preprocessing & Encoding, scaling, model-dependent transforms \\
\protect\hyperlink{chapter-10-training-inference}{\textbf{10}} &
Training \& Inference & Spine design, dataset generation \\
\protect\hyperlink{chapter-11-operations}{\textbf{11}} & Operations &
Monitoring, DMFs, lineage, drift detection \\
\protect\hyperlink{chapter-12-advanced-patterns}{\textbf{12}} & Advanced
Patterns & Streaming FeatureViews, multi-FS operations \\
\protect\hyperlink{chapter-13-migration-guide}{\textbf{13}} & Migration
Guide & Tecton, SageMaker, Vertex AI migration paths \\
\end{longtable}

\hypertarget{prerequisites-1}{%
\section{Prerequisites}\label{prerequisites-1}}

\hypertarget{required-snowflake-access}{%
\subsection{Required Snowflake Access}\label{required-snowflake-access}}

\begin{itemize}
\tightlist
\item
  Snowflake account (Standard Edition or higher)
\item
  Database privileges to create schemas, tables, dynamic tables, views,
  and tags
\item
  Warehouse access for compute
\item
  \texttt{snowflake-ml-python} package installed (available via
  pip/conda)
\end{itemize}

\hypertarget{required-privileges}{%
\subsection{Required Privileges}\label{required-privileges}}

Feature Store requires specific privileges. A role setup script is
available in the
\href{https://docs.snowflake.com/en/developer-guide/snowflake-ml/feature-store/rbac}{Snowflake
documentation}.

Minimum required privileges:

\begin{Shaded}
\begin{Highlighting}[]
\CommentTok{{-}{-} Create a role for Feature Store users}
\KeywordTok{CREATE} \KeywordTok{ROLE} \ControlFlowTok{IF} \KeywordTok{NOT} \KeywordTok{EXISTS}\NormalTok{ FEATURE\_STORE\_USER;}

\CommentTok{{-}{-} Grant database and schema access}
\KeywordTok{GRANT} \KeywordTok{USAGE} \KeywordTok{ON} \KeywordTok{DATABASE} \OperatorTok{\textless{}}\KeywordTok{database}\OperatorTok{\textgreater{}} \KeywordTok{TO} \KeywordTok{ROLE}\NormalTok{ FEATURE\_STORE\_USER;}
\KeywordTok{GRANT} \KeywordTok{CREATE} \KeywordTok{SCHEMA} \KeywordTok{ON} \KeywordTok{DATABASE} \OperatorTok{\textless{}}\KeywordTok{database}\OperatorTok{\textgreater{}} \KeywordTok{TO} \KeywordTok{ROLE}\NormalTok{ FEATURE\_STORE\_USER;}
\KeywordTok{GRANT} \KeywordTok{USAGE} \KeywordTok{ON}\NormalTok{ WAREHOUSE }\OperatorTok{\textless{}}\NormalTok{warehouse}\OperatorTok{\textgreater{}} \KeywordTok{TO} \KeywordTok{ROLE}\NormalTok{ FEATURE\_STORE\_USER;}

\CommentTok{{-}{-} For Feature Store schema operations}
\KeywordTok{GRANT} \KeywordTok{CREATE} \KeywordTok{TABLE} \KeywordTok{ON} \KeywordTok{SCHEMA} \OperatorTok{\textless{}}\KeywordTok{schema}\OperatorTok{\textgreater{}} \KeywordTok{TO} \KeywordTok{ROLE}\NormalTok{ FEATURE\_STORE\_USER;}
\KeywordTok{GRANT} \KeywordTok{CREATE} \KeywordTok{DYNAMIC} \KeywordTok{TABLE} \KeywordTok{ON} \KeywordTok{SCHEMA} \OperatorTok{\textless{}}\KeywordTok{schema}\OperatorTok{\textgreater{}} \KeywordTok{TO} \KeywordTok{ROLE}\NormalTok{ FEATURE\_STORE\_USER;}
\KeywordTok{GRANT} \KeywordTok{CREATE} \KeywordTok{VIEW} \KeywordTok{ON} \KeywordTok{SCHEMA} \OperatorTok{\textless{}}\KeywordTok{schema}\OperatorTok{\textgreater{}} \KeywordTok{TO} \KeywordTok{ROLE}\NormalTok{ FEATURE\_STORE\_USER;}
\KeywordTok{GRANT} \KeywordTok{CREATE}\NormalTok{ TAG }\KeywordTok{ON} \KeywordTok{SCHEMA} \OperatorTok{\textless{}}\KeywordTok{schema}\OperatorTok{\textgreater{}} \KeywordTok{TO} \KeywordTok{ROLE}\NormalTok{ FEATURE\_STORE\_USER;}
\end{Highlighting}
\end{Shaded}

\hypertarget{software-requirements}{%
\subsection{Software Requirements}\label{software-requirements}}

\begin{Shaded}
\begin{Highlighting}[]
\CommentTok{\# Minimum package versions for this guide}
\NormalTok{snowflake}\OperatorTok{{-}}\NormalTok{ml}\OperatorTok{{-}}\NormalTok{python }\OperatorTok{\textgreater{}=} \FloatTok{1.8.0}      \CommentTok{\# Core Feature Store functionality}
\NormalTok{snowflake}\OperatorTok{{-}}\NormalTok{snowpark}\OperatorTok{{-}}\NormalTok{python }\OperatorTok{\textgreater{}=} \FloatTok{1.25.0}  \CommentTok{\# DataFrame operations}

\CommentTok{\# Additional packages used in examples}
\NormalTok{pandas }\OperatorTok{\textgreater{}=} \FloatTok{2.0.0}
\NormalTok{numpy }\OperatorTok{\textgreater{}=} \FloatTok{1.23.5}
\NormalTok{scikit}\OperatorTok{{-}}\NormalTok{learn }\OperatorTok{\textgreater{}=} \FloatTok{1.5.0}
\end{Highlighting}
\end{Shaded}

\begin{tcolorbox}[enhanced jigsaw, title=\textcolor{quarto-callout-note-color}{\faInfo}\hspace{0.5em}{🆕 New Features in snowflake-ml 1.21+}, coltitle=black, left=2mm, colback=white, toptitle=1mm, opacityback=0, colbacktitle=quarto-callout-note-color!10!white, toprule=.15mm, colframe=quarto-callout-note-color-frame, leftrule=.75mm, opacitybacktitle=0.6, breakable, bottomtitle=1mm, arc=.35mm, bottomrule=.15mm, titlerule=0mm, rightrule=.15mm]

The \texttt{Feature} aggregation class
(Section~\ref{sec-aggregations-api}) and Temporal Aggregated
FeatureViews require version 1.21.0 or later. Earlier versions can use
the approaches shown in \textbf{?@sec-feature-pipelines} for temporal
features.

\end{tcolorbox}

\hypertarget{environment-setup}{%
\section{Environment Setup}\label{environment-setup}}

\hypertarget{session-configuration}{%
\subsection{Session Configuration}\label{session-configuration}}

The following code establishes a Snowpark session and verifies the
environment.

\begin{quote}
📁 \textbf{Full code:}
\href{_code/setup_session.py}{\texttt{\_code/setup\_session.py}}
\end{quote}

\begin{Shaded}
\begin{Highlighting}[]
\CommentTok{\# ==============================================================================}
\CommentTok{\# CONFIGURATION}
\CommentTok{\# ==============================================================================}

\CommentTok{\# Default configuration {-} override as needed}
\NormalTok{SOURCE\_DATABASE }\OperatorTok{=} \StringTok{\textquotesingle{}FEATURE\_STORE\_GUIDE\textquotesingle{}}
\NormalTok{SOURCE\_SCHEMA }\OperatorTok{=} \StringTok{\textquotesingle{}CLICKSTREAM\_RAW\textquotesingle{}}
\NormalTok{FS\_NAME }\OperatorTok{=} \StringTok{\textquotesingle{}GUIDE\_FEATURE\_STORE\textquotesingle{}}
\NormalTok{WAREHOUSE }\OperatorTok{=} \StringTok{\textquotesingle{}COMPUTE\_WH\textquotesingle{}}

\CommentTok{\# ==============================================================================}
\CommentTok{\# IMPORTS}
\CommentTok{\# ==============================================================================}

\CommentTok{\# Python standard library}
\ImportTok{import}\NormalTok{ json}
\ImportTok{from}\NormalTok{ datetime }\ImportTok{import}\NormalTok{ datetime, timedelta}
\ImportTok{from}\NormalTok{ decimal }\ImportTok{import}\NormalTok{ Decimal}

\CommentTok{\# Data processing}
\ImportTok{import}\NormalTok{ pandas }\ImportTok{as}\NormalTok{ pd}
\ImportTok{import}\NormalTok{ numpy }\ImportTok{as}\NormalTok{ np}

\CommentTok{\# Snowpark}
\ImportTok{from}\NormalTok{ snowflake.snowpark }\ImportTok{import}\NormalTok{ Session}
\ImportTok{from}\NormalTok{ snowflake.snowpark.version }\ImportTok{import}\NormalTok{ VERSION}
\ImportTok{from}\NormalTok{ snowflake.snowpark }\ImportTok{import}\NormalTok{ functions }\ImportTok{as}\NormalTok{ F}
\ImportTok{from}\NormalTok{ snowflake.snowpark }\ImportTok{import}\NormalTok{ types }\ImportTok{as}\NormalTok{ T}
\ImportTok{from}\NormalTok{ snowflake.snowpark.context }\ImportTok{import}\NormalTok{ get\_active\_session}

\CommentTok{\# Feature Store}
\ImportTok{from}\NormalTok{ snowflake.ml.feature\_store }\ImportTok{import}\NormalTok{ (}
\NormalTok{    FeatureStore, }
\NormalTok{    FeatureView, }
\NormalTok{    Entity,}
\NormalTok{    CreationMode}
\NormalTok{)}

\CommentTok{\# ==============================================================================}
\CommentTok{\# SESSION CREATION}
\CommentTok{\# ==============================================================================}

\CommentTok{\# Option 1: Running in Snowflake Notebook (preferred)}
\ControlFlowTok{try}\NormalTok{:}
\NormalTok{    session }\OperatorTok{=}\NormalTok{ get\_active\_session()}
\ControlFlowTok{except}\NormalTok{:}
    \CommentTok{\# Option 2: Running locally with connection config}
    \ImportTok{from}\NormalTok{ snowflake.ml.utils }\ImportTok{import}\NormalTok{ connection\_params}
\NormalTok{    session }\OperatorTok{=}\NormalTok{ Session.builder.configs(connection\_params.SnowflakeLoginOptions()).create()}

\CommentTok{\# Enable SQL simplifier for cleaner generated SQL}
\NormalTok{session.sql\_simplifier\_enabled }\OperatorTok{=} \VariableTok{True}
\end{Highlighting}
\end{Shaded}

\hypertarget{verify-environment}{%
\subsection{Verify Environment}\label{verify-environment}}

\begin{quote}
📁 \textbf{Full code:}
\href{_code/verify_environment.py}{\texttt{\_code/verify\_environment.py}}
\end{quote}

\begin{Shaded}
\begin{Highlighting}[]
\CommentTok{\# Capture environment details}
\NormalTok{snowflake\_environment }\OperatorTok{=}\NormalTok{ session.sql(}\StringTok{\textquotesingle{}SELECT current\_user(), current\_version()\textquotesingle{}}\NormalTok{).collect()}
\NormalTok{snowpark\_version }\OperatorTok{=}\NormalTok{ VERSION}
\NormalTok{session\_role }\OperatorTok{=}\NormalTok{ session.get\_current\_role().replace(}\StringTok{\textquotesingle{}"\textquotesingle{}}\NormalTok{, }\StringTok{""}\NormalTok{)}
\NormalTok{session\_database }\OperatorTok{=}\NormalTok{ session.get\_current\_database().replace(}\StringTok{\textquotesingle{}"\textquotesingle{}}\NormalTok{, }\StringTok{""}\NormalTok{)}
\NormalTok{session\_schema }\OperatorTok{=}\NormalTok{ session.get\_current\_schema().replace(}\StringTok{\textquotesingle{}"\textquotesingle{}}\NormalTok{, }\StringTok{""}\NormalTok{)}
\NormalTok{session\_warehouse }\OperatorTok{=}\NormalTok{ session.get\_current\_warehouse().replace(}\StringTok{\textquotesingle{}"\textquotesingle{}}\NormalTok{, }\StringTok{""}\NormalTok{)}

\CommentTok{\# Check warehouse status}
\NormalTok{wh\_status }\OperatorTok{=}\NormalTok{ session.sql(}\SpecialStringTok{f"SHOW WAREHOUSES LIKE \textquotesingle{}}\SpecialCharTok{\{}\NormalTok{session\_warehouse}\SpecialCharTok{\}}\SpecialStringTok{\textquotesingle{}"}\NormalTok{).collect()[}\DecValTok{0}\NormalTok{]}

\BuiltInTok{print}\NormalTok{(}\StringTok{\textquotesingle{}=\textquotesingle{}} \OperatorTok{*} \DecValTok{70}\NormalTok{)}
\BuiltInTok{print}\NormalTok{(}\StringTok{\textquotesingle{}CONNECTION ESTABLISHED\textquotesingle{}}\NormalTok{)}
\BuiltInTok{print}\NormalTok{(}\StringTok{\textquotesingle{}=\textquotesingle{}} \OperatorTok{*} \DecValTok{70}\NormalTok{)}
\BuiltInTok{print}\NormalTok{(}\SpecialStringTok{f\textquotesingle{}Account                : }\SpecialCharTok{\{}\NormalTok{session}\SpecialCharTok{.}\NormalTok{sql(}\StringTok{"SELECT current\_account()"}\NormalTok{)}\SpecialCharTok{.}\NormalTok{collect()[}\DecValTok{0}\NormalTok{][}\DecValTok{0}\NormalTok{]}\SpecialCharTok{\}}\SpecialStringTok{\textquotesingle{}}\NormalTok{)}
\BuiltInTok{print}\NormalTok{(}\SpecialStringTok{f\textquotesingle{}User                   : }\SpecialCharTok{\{}\NormalTok{snowflake\_environment[}\DecValTok{0}\NormalTok{][}\DecValTok{0}\NormalTok{]}\SpecialCharTok{\}}\SpecialStringTok{\textquotesingle{}}\NormalTok{)}
\BuiltInTok{print}\NormalTok{(}\SpecialStringTok{f\textquotesingle{}Role                   : }\SpecialCharTok{\{}\NormalTok{session\_role}\SpecialCharTok{\}}\SpecialStringTok{\textquotesingle{}}\NormalTok{)}
\BuiltInTok{print}\NormalTok{(}\SpecialStringTok{f\textquotesingle{}Database               : }\SpecialCharTok{\{}\NormalTok{session\_database}\SpecialCharTok{\}}\SpecialStringTok{\textquotesingle{}}\NormalTok{)}
\BuiltInTok{print}\NormalTok{(}\SpecialStringTok{f\textquotesingle{}Schema                 : }\SpecialCharTok{\{}\NormalTok{session\_schema}\SpecialCharTok{\}}\SpecialStringTok{\textquotesingle{}}\NormalTok{)}
\BuiltInTok{print}\NormalTok{(}\SpecialStringTok{f\textquotesingle{}Warehouse              : }\SpecialCharTok{\{}\NormalTok{session\_warehouse}\SpecialCharTok{\}}\SpecialStringTok{\textquotesingle{}}\NormalTok{)}
\BuiltInTok{print}\NormalTok{(}\SpecialStringTok{f\textquotesingle{}Warehouse Size         : }\SpecialCharTok{\{}\NormalTok{wh\_status[}\StringTok{"size"}\NormalTok{]}\SpecialCharTok{.}\NormalTok{upper()}\SpecialCharTok{\}}\SpecialStringTok{\textquotesingle{}}\NormalTok{)}
\BuiltInTok{print}\NormalTok{(}\SpecialStringTok{f\textquotesingle{}Warehouse State        : }\SpecialCharTok{\{}\NormalTok{wh\_status[}\StringTok{"state"}\NormalTok{]}\SpecialCharTok{\}}\SpecialStringTok{\textquotesingle{}}\NormalTok{)}
\BuiltInTok{print}\NormalTok{(}\SpecialStringTok{f\textquotesingle{}Snowflake Version      : }\SpecialCharTok{\{}\NormalTok{snowflake\_environment[}\DecValTok{0}\NormalTok{][}\DecValTok{1}\NormalTok{]}\SpecialCharTok{\}}\SpecialStringTok{\textquotesingle{}}\NormalTok{)}
\BuiltInTok{print}\NormalTok{(}\SpecialStringTok{f\textquotesingle{}Snowpark Version       : }\SpecialCharTok{\{}\NormalTok{snowpark\_version[}\DecValTok{0}\NormalTok{]}\SpecialCharTok{\}}\SpecialStringTok{.}\SpecialCharTok{\{}\NormalTok{snowpark\_version[}\DecValTok{1}\NormalTok{]}\SpecialCharTok{\}}\SpecialStringTok{.}\SpecialCharTok{\{}\NormalTok{snowpark\_version[}\DecValTok{2}\NormalTok{]}\SpecialCharTok{\}}\SpecialStringTok{\textquotesingle{}}\NormalTok{)}
\BuiltInTok{print}\NormalTok{(}\StringTok{\textquotesingle{}=\textquotesingle{}} \OperatorTok{*} \DecValTok{70}\NormalTok{)}
\end{Highlighting}
\end{Shaded}

Expected output:

\begin{verbatim}
======================================================================
CONNECTION ESTABLISHED
======================================================================
Account                : YOUR_ACCOUNT
User                   : YOUR_USER
Role                   : SYSADMIN
Database               : FEATURE_STORE_GUIDE
Schema                 : CLICKSTREAM_RAW
Warehouse              : COMPUTE_WH
Warehouse Size         : X-SMALL
Warehouse State        : STARTED
Snowflake Version      : 8.x.x
Snowpark Version       : 1.25.0
======================================================================
\end{verbatim}

\hypertarget{create-feature-store}{%
\subsection{Create Feature Store}\label{create-feature-store}}

\begin{quote}
📁 \textbf{Full code:}
\href{_code/create_feature_store.py}{\texttt{\_code/create\_feature\_store.py}}
\end{quote}

\begin{Shaded}
\begin{Highlighting}[]
\CommentTok{\# Create working schema for sample data}
\NormalTok{session.sql(}\SpecialStringTok{f\textquotesingle{}CREATE SCHEMA IF NOT EXISTS }\SpecialCharTok{\{}\NormalTok{SOURCE\_DATABASE}\SpecialCharTok{\}}\SpecialStringTok{.}\SpecialCharTok{\{}\NormalTok{SOURCE\_SCHEMA}\SpecialCharTok{\}}\SpecialStringTok{\textquotesingle{}}\NormalTok{).collect()}

\CommentTok{\# Initialize Feature Store}
\CommentTok{\# This creates a schema with Feature Store metadata tags}
\NormalTok{fs }\OperatorTok{=}\NormalTok{ FeatureStore(}
\NormalTok{    session}\OperatorTok{=}\NormalTok{session,}
\NormalTok{    database}\OperatorTok{=}\NormalTok{SOURCE\_DATABASE,}
\NormalTok{    name}\OperatorTok{=}\NormalTok{FS\_NAME,}
\NormalTok{    default\_warehouse}\OperatorTok{=}\NormalTok{WAREHOUSE,}
\NormalTok{    creation\_mode}\OperatorTok{=}\NormalTok{CreationMode.CREATE\_IF\_NOT\_EXIST,}
\NormalTok{)}

\BuiltInTok{print}\NormalTok{(}\SpecialStringTok{f\textquotesingle{}Feature Store initialized: }\SpecialCharTok{\{}\NormalTok{SOURCE\_DATABASE}\SpecialCharTok{\}}\SpecialStringTok{.}\SpecialCharTok{\{}\NormalTok{FS\_NAME}\SpecialCharTok{\}}\SpecialStringTok{\textquotesingle{}}\NormalTok{)}
\end{Highlighting}
\end{Shaded}

\hypertarget{sample-data}{%
\section{Sample Data}\label{sample-data}}

This guide uses a synthetic clickstream dataset designed to demonstrate
all Feature Store capabilities, including:

\begin{itemize}
\tightlist
\item
  \textbf{Temporal patterns}: Event timestamps with realistic
  distributions
\item
  \textbf{Entity hierarchies}: Visitor → User → Household relationships
\item
  \textbf{Composite keys}: Product-Supplier junction table (M:N)
\item
  \textbf{Semi-structured data}: VARIANT, ARRAY, OBJECT columns
\end{itemize}

See \protect\hyperlink{appendix-a-sample-data}{Appendix A: Clickstream
Data Model} for the full data model specification.

\hypertarget{quick-data-overview}{%
\subsection{Quick Data Overview}\label{quick-data-overview}}

The clickstream dataset includes these core tables:

\begin{longtable}[]{@{}
  >{\raggedright\arraybackslash}p{(\columnwidth - 6\tabcolsep) * \real{0.1429}}
  >{\raggedright\arraybackslash}p{(\columnwidth - 6\tabcolsep) * \real{0.2653}}
  >{\raggedright\arraybackslash}p{(\columnwidth - 6\tabcolsep) * \real{0.2449}}
  >{\raggedright\arraybackslash}p{(\columnwidth - 6\tabcolsep) * \real{0.3469}}@{}}
\toprule\noalign{}
\begin{minipage}[b]{\linewidth}\raggedright
Table
\end{minipage} & \begin{minipage}[b]{\linewidth}\raggedright
Description
\end{minipage} & \begin{minipage}[b]{\linewidth}\raggedright
Entity Key
\end{minipage} & \begin{minipage}[b]{\linewidth}\raggedright
Temporal Column
\end{minipage} \\
\midrule\noalign{}
\endhead
\bottomrule\noalign{}
\endlastfoot
\texttt{VISITORS} & Anonymous visitors & \texttt{VISITOR\_ID} &
\texttt{FIRST\_SEEN\_TS}, \texttt{LAST\_SEEN\_TS} \\
\texttt{USERS} & Identified users & \texttt{USER\_ID} &
\texttt{CREATED\_TS}, \texttt{UPDATED\_TS} \\
\texttt{SESSIONS} & User sessions & \texttt{SESSION\_ID} &
\texttt{SESSION\_START\_TS}, \texttt{SESSION\_END\_TS} \\
\texttt{EVENTS} & Clickstream events & \texttt{EVENT\_ID} &
\texttt{EVENT\_TS} \\
\texttt{ORDERS} & Purchase orders & \texttt{ORDER\_ID} &
\texttt{ORDER\_TS} \\
\texttt{PRODUCTS} & Product catalog & \texttt{PRODUCT\_ID} &
\texttt{CREATED\_TS} \\
\texttt{PRODUCT\_SUPPLIER} & Supplier relationships &
\texttt{(PRODUCT\_ID,\ SUPPLIER\_ID)} & \texttt{VALID\_FROM\_TS},
\texttt{VALID\_TO\_TS} \\
\end{longtable}

\hypertarget{inline-snippet-pattern}{%
\subsection{Inline Snippet Pattern}\label{inline-snippet-pattern}}

For simple examples, we use self-contained data that you can copy-paste
directly:

\begin{Shaded}
\begin{Highlighting}[]
\CommentTok{\# Example: Create sample events for quick testing}
\ImportTok{from}\NormalTok{ snowflake.snowpark }\ImportTok{import}\NormalTok{ Session, Row}

\CommentTok{\# Assumes session is already created}
\NormalTok{sample\_events }\OperatorTok{=}\NormalTok{ session.create\_dataframe([}
\NormalTok{    Row(USER\_ID}\OperatorTok{=}\StringTok{\textquotesingle{}usr\_001\textquotesingle{}}\NormalTok{, EVENT\_TS}\OperatorTok{=}\StringTok{\textquotesingle{}2025{-}01{-}01 10:00:00\textquotesingle{}}\NormalTok{, EVENT\_TYPE}\OperatorTok{=}\StringTok{\textquotesingle{}page\_view\textquotesingle{}}\NormalTok{, VALUE}\OperatorTok{=}\FloatTok{1.0}\NormalTok{),}
\NormalTok{    Row(USER\_ID}\OperatorTok{=}\StringTok{\textquotesingle{}usr\_001\textquotesingle{}}\NormalTok{, EVENT\_TS}\OperatorTok{=}\StringTok{\textquotesingle{}2025{-}01{-}01 10:05:00\textquotesingle{}}\NormalTok{, EVENT\_TYPE}\OperatorTok{=}\StringTok{\textquotesingle{}click\textquotesingle{}}\NormalTok{, VALUE}\OperatorTok{=}\FloatTok{1.0}\NormalTok{),}
\NormalTok{    Row(USER\_ID}\OperatorTok{=}\StringTok{\textquotesingle{}usr\_001\textquotesingle{}}\NormalTok{, EVENT\_TS}\OperatorTok{=}\StringTok{\textquotesingle{}2025{-}01{-}01 10:10:00\textquotesingle{}}\NormalTok{, EVENT\_TYPE}\OperatorTok{=}\StringTok{\textquotesingle{}purchase\textquotesingle{}}\NormalTok{, VALUE}\OperatorTok{=}\FloatTok{99.0}\NormalTok{),}
\NormalTok{    Row(USER\_ID}\OperatorTok{=}\StringTok{\textquotesingle{}usr\_002\textquotesingle{}}\NormalTok{, EVENT\_TS}\OperatorTok{=}\StringTok{\textquotesingle{}2025{-}01{-}01 11:00:00\textquotesingle{}}\NormalTok{, EVENT\_TYPE}\OperatorTok{=}\StringTok{\textquotesingle{}page\_view\textquotesingle{}}\NormalTok{, VALUE}\OperatorTok{=}\FloatTok{1.0}\NormalTok{),}
\NormalTok{    Row(USER\_ID}\OperatorTok{=}\StringTok{\textquotesingle{}usr\_002\textquotesingle{}}\NormalTok{, EVENT\_TS}\OperatorTok{=}\StringTok{\textquotesingle{}2025{-}01{-}01 11:02:00\textquotesingle{}}\NormalTok{, EVENT\_TYPE}\OperatorTok{=}\StringTok{\textquotesingle{}page\_view\textquotesingle{}}\NormalTok{, VALUE}\OperatorTok{=}\FloatTok{1.0}\NormalTok{),}
\NormalTok{])}

\NormalTok{sample\_events.write.save\_as\_table(}\StringTok{"SAMPLE\_EVENTS"}\NormalTok{, mode}\OperatorTok{=}\StringTok{"overwrite"}\NormalTok{)}
\end{Highlighting}
\end{Shaded}

\begin{tcolorbox}[enhanced jigsaw, title=\textcolor{quarto-callout-tip-color}{\faLightbulb}\hspace{0.5em}{Full Code Available}, coltitle=black, left=2mm, colback=white, toptitle=1mm, opacityback=0, colbacktitle=quarto-callout-tip-color!10!white, toprule=.15mm, colframe=quarto-callout-tip-color-frame, leftrule=.75mm, opacitybacktitle=0.6, breakable, bottomtitle=1mm, arc=.35mm, bottomrule=.15mm, titlerule=0mm, rightrule=.15mm]

All code examples in this guide are available in the \texttt{\_code/}
directories. You can also click the ``View source'' icon on any code
block to see the full file.

\end{tcolorbox}

\hypertarget{guide-conventions}{%
\section{Guide Conventions}\label{guide-conventions}}

\hypertarget{code-examples-1}{%
\subsection{Code Examples}\label{code-examples-1}}

All code examples follow these conventions:

\begin{itemize}
\tightlist
\item
  \textbf{SQL}: UPPERCASE keywords, lowercase identifiers
\item
  \textbf{Python}: PEP 8 style, type hints where helpful
\item
  \textbf{Naming}: \texttt{SCREAMING\_SNAKE\_CASE} for Snowflake
  objects, \texttt{snake\_case} for Python
\end{itemize}

\hypertarget{sec-naming-conventions}{%
\subsection{Feature Naming Conventions}\label{sec-naming-conventions}}

\begin{longtable}[]{@{}lll@{}}
\toprule\noalign{}
Suffix & Meaning & Example \\
\midrule\noalign{}
\endhead
\bottomrule\noalign{}
\endlastfoot
\texttt{\_TS} & Timestamp & \texttt{ORDER\_TS}, \texttt{CREATED\_TS} \\
\texttt{\_DT} & Date & \texttt{ORDER\_DT} \\
\texttt{\_CNT} & Count & \texttt{SESSION\_CNT}, \texttt{ORDER\_CNT} \\
\texttt{\_DCNT} & Distinct count & \texttt{PRODUCT\_VIEW\_DCNT} \\
\texttt{\_SUM} & Sum/total & \texttt{REVENUE\_SUM},
\texttt{QUANTITY\_SUM} \\
\texttt{\_AVG} & Average & \texttt{ORDER\_VALUE\_AVG} \\
\texttt{\_AMT} & Amount (currency) & \texttt{TOTAL\_AMT},
\texttt{DISCOUNT\_AMT} \\
\texttt{IS\_} & Boolean flag & \texttt{IS\_CONVERTED},
\texttt{IS\_CURRENT} \\
\end{longtable}

\hypertarget{version-callouts}{%
\subsection{Version Callouts}\label{version-callouts}}

New features are highlighted with version callouts:

\begin{tcolorbox}[enhanced jigsaw, title=\textcolor{quarto-callout-note-color}{\faInfo}\hspace{0.5em}{🆕 New in snowflake-ml X.X.X}, coltitle=black, left=2mm, colback=white, toptitle=1mm, opacityback=0, colbacktitle=quarto-callout-note-color!10!white, toprule=.15mm, colframe=quarto-callout-note-color-frame, leftrule=.75mm, opacitybacktitle=0.6, breakable, bottomtitle=1mm, arc=.35mm, bottomrule=.15mm, titlerule=0mm, rightrule=.15mm]

Description of the new feature and how it improves on previous
approaches.

\end{tcolorbox}

\begin{center}\rule{0.5\linewidth}{0.5pt}\end{center}

\hypertarget{next-steps}{%
\section{Next Steps}\label{next-steps}}

Continue to \protect\hyperlink{core-concepts}{Chapter 01: Concepts} to
learn about the core Feature Store objects and terminology.

\hypertarget{core-concepts}{%
\chapter{Core Concepts}\label{core-concepts}}

Entities, FeatureViews, Features, and the Spine

\hfill\break

\hypertarget{overview-1}{%
\section{Overview}\label{overview-1}}

This chapter introduces the core concepts and terminology of Snowflake
Feature Store. Understanding these fundamentals is essential before
diving into implementation patterns.

\hypertarget{learning-objectives}{%
\section{Learning Objectives}\label{learning-objectives}}

After completing this chapter, you will be able to:

\begin{itemize}
\tightlist
\item
  Define what a Feature Store is and articulate its value in ML systems
\item
  Identify and describe the key components: Entities, FeatureViews,
  Features, and Spines
\item
  Understand how Feature Store fits into the broader ML lifecycle
\item
  Apply the correct transformation taxonomy (MIT/MDT/ODT) to your use
  cases
\item
  Map Snowflake terminology to industry-standard terms
\end{itemize}

\begin{center}\rule{0.5\linewidth}{0.5pt}\end{center}

\hypertarget{what-is-a-feature-store}{%
\section{What is a Feature Store?}\label{what-is-a-feature-store}}

A \textbf{Feature Store} is a centralized repository for storing,
managing, and serving features for machine learning. It acts as the data
management layer between raw data sources and ML models.

\begin{verbatim}
┌─────────────────────────────────────────────────────────────────────────────┐
│                           ML DATA INFRASTRUCTURE                            │
├─────────────────────────────────────────────────────────────────────────────┤
│                                                                             │
│   ┌──────────────┐     ┌────────────────────┐     ┌──────────────────────┐  │
│   │ Data Sources │     │   Feature Store    │     │    ML Consumers      │  │
│   │              │     │                    │     │                      │  │
│   │ • Tables     │────▶│ • Feature Pipelines│────▶│ • Model Training     │  │
│   │ • Streams    │     │ • Feature Storage  │     │ • Batch Inference    │  │
│   │ • External   │     │ • Feature Serving  │     │ • Real-time Serving  │  │
│   │ • APIs       │     │ • Metadata & Gov   │     │ • Analytics          │  │
│   └──────────────┘     └────────────────────┘     └──────────────────────┘  │
│                                                                             │
└─────────────────────────────────────────────────────────────────────────────┘
\end{verbatim}

\hypertarget{features-the-building-blocks-of-ml}{%
\subsection{Features: The Building Blocks of
ML}\label{features-the-building-blocks-of-ml}}

A \textbf{feature} is a measurable property or characteristic used as
input to an ML model. Features can be:

\begin{longtable}[]{@{}
  >{\raggedright\arraybackslash}p{(\columnwidth - 4\tabcolsep) * \real{0.3889}}
  >{\raggedright\arraybackslash}p{(\columnwidth - 4\tabcolsep) * \real{0.3611}}
  >{\raggedright\arraybackslash}p{(\columnwidth - 4\tabcolsep) * \real{0.2500}}@{}}
\toprule\noalign{}
\begin{minipage}[b]{\linewidth}\raggedright
Feature Type
\end{minipage} & \begin{minipage}[b]{\linewidth}\raggedright
Description
\end{minipage} & \begin{minipage}[b]{\linewidth}\raggedright
Example
\end{minipage} \\
\midrule\noalign{}
\endhead
\bottomrule\noalign{}
\endlastfoot
\textbf{Raw} & Direct attributes from source data & \texttt{user\_age},
\texttt{product\_price} \\
\textbf{Derived} & Calculated from raw features &
\texttt{price\_per\_unit}, \texttt{age\_bucket} \\
\textbf{Aggregated} & Computed over time windows &
\texttt{total\_orders\_30d}, \texttt{avg\_session\_duration\_7d} \\
\textbf{Encoded} & Transformed for model consumption &
\texttt{category\_one\_hot}, \texttt{amount\_scaled} \\
\end{longtable}

\begin{center}\rule{0.5\linewidth}{0.5pt}\end{center}

\hypertarget{why-use-a-feature-store}{%
\section{Why Use a Feature Store?}\label{why-use-a-feature-store}}

Feature Stores solve critical challenges in production ML systems:

\hypertarget{feature-reusability}{%
\subsection{1. Feature Reusability}\label{feature-reusability}}

Without a Feature Store, teams often recreate the same features
independently, leading to:

\begin{itemize}
\tightlist
\item
  Duplicated effort across data scientists
\item
  Inconsistent feature definitions
\item
  Wasted compute resources
\end{itemize}

\textbf{With Feature Store}: Features are defined once, stored
centrally, and reused across models and teams.

\hypertarget{training-serving-skew}{%
\subsection{2. Training-Serving Skew}\label{training-serving-skew}}

One of the most common causes of ML model degradation is when features
computed during training differ from those used during inference.

\textbf{With Feature Store}: The same feature definitions serve both
training and inference, eliminating skew.

\hypertarget{point-in-time-correctness}{%
\subsection{3. Point-in-Time
Correctness}\label{point-in-time-correctness}}

Using future data to train models (data leakage) produces artificially
good training metrics but poor production performance.

\textbf{With Feature Store}: Built-in temporal joins ensure features are
computed using only data available at the prediction time, preventing
data leakage.

\hypertarget{feature-discovery}{%
\subsection{4. Feature Discovery}\label{feature-discovery}}

As organizations scale ML, finding existing features becomes
challenging.

\textbf{With Feature Store}: Searchable catalog of features with
metadata, lineage, and documentation.

\hypertarget{governance-compliance}{%
\subsection{5. Governance \& Compliance}\label{governance-compliance}}

Tracking where features come from and how they're used is essential for
regulatory compliance.

\textbf{With Feature Store}: Full lineage from source data through
feature transformations to model predictions.

\begin{center}\rule{0.5\linewidth}{0.5pt}\end{center}

\hypertarget{key-components}{%
\section{Key Components}\label{key-components}}

Snowflake Feature Store consists of four primary components:

\hypertarget{sec-entity}{%
\subsection{Entity}\label{sec-entity}}

An \textbf{Entity} represents a business object that features describe.
It defines the join keys used to retrieve features.

\begin{quote}
📁 \textbf{Full code:}
\href{_code/entity_examples.py}{\texttt{\_code/entity\_examples.py}}
\end{quote}

\begin{Shaded}
\begin{Highlighting}[]
\ImportTok{from}\NormalTok{ snowflake.ml.feature\_store }\ImportTok{import}\NormalTok{ Entity}

\CommentTok{\# Simple entity with single key}
\NormalTok{user\_entity }\OperatorTok{=}\NormalTok{ Entity(}
\NormalTok{    name}\OperatorTok{=}\StringTok{"USER"}\NormalTok{,}
\NormalTok{    join\_keys}\OperatorTok{=}\NormalTok{[}\StringTok{"USER\_ID"}\NormalTok{],}
\NormalTok{    desc}\OperatorTok{=}\StringTok{"Registered user in the system"}
\NormalTok{)}

\CommentTok{\# Compound entity with multiple keys}
\NormalTok{order\_item\_entity }\OperatorTok{=}\NormalTok{ Entity(}
\NormalTok{    name}\OperatorTok{=}\StringTok{"ORDER\_ITEM"}\NormalTok{,}
\NormalTok{    join\_keys}\OperatorTok{=}\NormalTok{[}\StringTok{"ORDER\_ID"}\NormalTok{, }\StringTok{"LINE\_NUMBER"}\NormalTok{],}
\NormalTok{    desc}\OperatorTok{=}\StringTok{"Individual line item within an order"}
\NormalTok{)}
\end{Highlighting}
\end{Shaded}

\textbf{Key Concepts}:

\begin{longtable}[]{@{}
  >{\raggedright\arraybackslash}p{(\columnwidth - 2\tabcolsep) * \real{0.4091}}
  >{\raggedright\arraybackslash}p{(\columnwidth - 2\tabcolsep) * \real{0.5909}}@{}}
\toprule\noalign{}
\begin{minipage}[b]{\linewidth}\raggedright
Concept
\end{minipage} & \begin{minipage}[b]{\linewidth}\raggedright
Description
\end{minipage} \\
\midrule\noalign{}
\endhead
\bottomrule\noalign{}
\endlastfoot
\texttt{name} & Unique identifier for the entity within the Feature
Store \\
\texttt{join\_keys} & Column(s) that uniquely identify an instance of
the entity \\
\textbf{Simple Key} & Single column identifier (e.g.,
\texttt{USER\_ID}) \\
\textbf{Compound Key} & Multiple columns together form the identifier
(e.g., \texttt{ORDER\_ID} + \texttt{LINE\_NUMBER}) \\
\end{longtable}

Entities enable features from different FeatureViews to be joined
together when generating training data or serving predictions.

\begin{tcolorbox}[enhanced jigsaw, title=\textcolor{quarto-callout-note-color}{\faInfo}\hspace{0.5em}{📖 Deep Dive}, coltitle=black, left=2mm, colback=white, toptitle=1mm, opacityback=0, colbacktitle=quarto-callout-note-color!10!white, toprule=.15mm, colframe=quarto-callout-note-color-frame, leftrule=.75mm, opacitybacktitle=0.6, breakable, bottomtitle=1mm, arc=.35mm, bottomrule=.15mm, titlerule=0mm, rightrule=.15mm]

See \protect\hyperlink{chapter-03-entities-hierarchies}{Chapter 03:
Entities \& Hierarchies} for entity design patterns, hierarchies, and
compound key strategies.

\end{tcolorbox}

\begin{center}\rule{0.5\linewidth}{0.5pt}\end{center}

\hypertarget{sec-featureview}{%
\subsection{FeatureView}\label{sec-featureview}}

A \textbf{FeatureView} is a collection of related features computed from
a source dataframe. It defines:

\begin{itemize}
\tightlist
\item
  Which entity the features belong to
\item
  How features are computed (the transformation logic)
\item
  How features are materialized (Dynamic Table vs View)
\end{itemize}

\begin{quote}
📁 \textbf{Full code:}
\href{_code/featureview_examples.py}{\texttt{\_code/featureview\_examples.py}}
\end{quote}

\begin{Shaded}
\begin{Highlighting}[]
\ImportTok{from}\NormalTok{ snowflake.ml.feature\_store }\ImportTok{import}\NormalTok{ FeatureView}

\CommentTok{\# Create a FeatureView from a Snowpark DataFrame}
\NormalTok{user\_features\_fv }\OperatorTok{=}\NormalTok{ FeatureView(}
\NormalTok{    name}\OperatorTok{=}\StringTok{"USER\_PURCHASE\_FEATURES"}\NormalTok{,}
\NormalTok{    entities}\OperatorTok{=}\NormalTok{[user\_entity],}
\NormalTok{    feature\_df}\OperatorTok{=}\NormalTok{user\_purchase\_df,  }\CommentTok{\# Snowpark DataFrame with feature logic}
\NormalTok{    timestamp\_col}\OperatorTok{=}\StringTok{"UPDATED\_TS"}\NormalTok{,   }\CommentTok{\# For point{-}in{-}time correctness}
\NormalTok{    refresh\_freq}\OperatorTok{=}\StringTok{"1 hour"}\NormalTok{,        }\CommentTok{\# Dynamic Table refresh (omit for View)}
\NormalTok{    desc}\OperatorTok{=}\StringTok{"User purchase behavior features"}
\NormalTok{)}

\CommentTok{\# Register in Feature Store}
\NormalTok{user\_features\_fv }\OperatorTok{=}\NormalTok{ fs.register\_feature\_view(}
\NormalTok{    feature\_view}\OperatorTok{=}\NormalTok{user\_features\_fv,}
\NormalTok{    version}\OperatorTok{=}\StringTok{"V1"}\NormalTok{,}
\NormalTok{    block}\OperatorTok{=}\VariableTok{True}  \CommentTok{\# Wait for initial materialization}
\NormalTok{)}
\end{Highlighting}
\end{Shaded}

\textbf{Materialization Options}:

\begin{longtable}[]{@{}
  >{\raggedright\arraybackslash}p{(\columnwidth - 6\tabcolsep) * \real{0.1333}}
  >{\raggedright\arraybackslash}p{(\columnwidth - 6\tabcolsep) * \real{0.3111}}
  >{\raggedright\arraybackslash}p{(\columnwidth - 6\tabcolsep) * \real{0.2222}}
  >{\raggedright\arraybackslash}p{(\columnwidth - 6\tabcolsep) * \real{0.3333}}@{}}
\toprule\noalign{}
\begin{minipage}[b]{\linewidth}\raggedright
Type
\end{minipage} & \begin{minipage}[b]{\linewidth}\raggedright
Created When
\end{minipage} & \begin{minipage}[b]{\linewidth}\raggedright
Use Case
\end{minipage} & \begin{minipage}[b]{\linewidth}\raggedright
Compute Model
\end{minipage} \\
\midrule\noalign{}
\endhead
\bottomrule\noalign{}
\endlastfoot
\textbf{Dynamic Table} & \texttt{refresh\_freq} is specified &
Pre-computed features, automatic refresh & Snowflake manages refresh \\
\textbf{View} & \texttt{refresh\_freq} is omitted & Query-time
computation, always fresh & Compute on each query \\
\end{longtable}

\textbf{Key Concepts}:

\begin{longtable}[]{@{}
  >{\raggedright\arraybackslash}p{(\columnwidth - 2\tabcolsep) * \real{0.4091}}
  >{\raggedright\arraybackslash}p{(\columnwidth - 2\tabcolsep) * \real{0.5909}}@{}}
\toprule\noalign{}
\begin{minipage}[b]{\linewidth}\raggedright
Concept
\end{minipage} & \begin{minipage}[b]{\linewidth}\raggedright
Description
\end{minipage} \\
\midrule\noalign{}
\endhead
\bottomrule\noalign{}
\endlastfoot
\texttt{name} & Unique identifier for the FeatureView \\
\texttt{entities} & List of entities this FeatureView provides features
for \\
\texttt{feature\_df} & Snowpark DataFrame defining the feature
transformations \\
\texttt{timestamp\_col} & Column used for point-in-time feature
retrieval \\
\texttt{refresh\_freq} & How often to refresh (e.g., \texttt{"1\ hour"},
\texttt{"1\ day"}) \\
\texttt{version} & Version string for managing FeatureView evolution \\
\end{longtable}

\begin{tcolorbox}[enhanced jigsaw, title=\textcolor{quarto-callout-note-color}{\faInfo}\hspace{0.5em}{📖 Deep Dive}, coltitle=black, left=2mm, colback=white, toptitle=1mm, opacityback=0, colbacktitle=quarto-callout-note-color!10!white, toprule=.15mm, colframe=quarto-callout-note-color-frame, leftrule=.75mm, opacitybacktitle=0.6, breakable, bottomtitle=1mm, arc=.35mm, bottomrule=.15mm, titlerule=0mm, rightrule=.15mm]

See \protect\hyperlink{chapter-04-feature-views}{Chapter 04: Feature
Views} for detailed coverage of FeatureView types, versioning, and
lifecycle management.

\end{tcolorbox}

\begin{center}\rule{0.5\linewidth}{0.5pt}\end{center}

\hypertarget{sec-feature}{%
\subsection{Feature}\label{sec-feature}}

A \textbf{Feature} is an individual column within a FeatureView.
Features are defined through the \texttt{feature\_df} parameter---a
Snowpark DataFrame that specifies the transformation logic.

\hypertarget{method-1-dataframe-api-most-common}{%
\subsubsection{Method 1: DataFrame API (Most
Common)}\label{method-1-dataframe-api-most-common}}

Use Snowpark DataFrame operations to define features declaratively:

\begin{quote}
📁 \textbf{Full code:}
\href{_code/feature_dataframe_api.py}{\texttt{\_code/feature\_dataframe\_api.py}}
\end{quote}

\begin{Shaded}
\begin{Highlighting}[]
\ImportTok{from}\NormalTok{ snowflake.snowpark }\ImportTok{import}\NormalTok{ functions }\ImportTok{as}\NormalTok{ F}

\CommentTok{\# Define features using Snowpark DataFrame transformations}
\NormalTok{user\_purchase\_df }\OperatorTok{=}\NormalTok{ (}
\NormalTok{    session.table(}\StringTok{"ORDERS"}\NormalTok{)}
\NormalTok{    .group\_by(}\StringTok{"USER\_ID"}\NormalTok{)}
\NormalTok{    .agg(}
\NormalTok{        F.}\BuiltInTok{sum}\NormalTok{(}\StringTok{"AMOUNT"}\NormalTok{).alias(}\StringTok{"TOTAL\_SPEND"}\NormalTok{),}
\NormalTok{        F.count(}\StringTok{"ORDER\_ID"}\NormalTok{).alias(}\StringTok{"ORDER\_CNT"}\NormalTok{),}
\NormalTok{        F.avg(}\StringTok{"AMOUNT"}\NormalTok{).alias(}\StringTok{"AVG\_ORDER\_AMT"}\NormalTok{),}
\NormalTok{        F.}\BuiltInTok{max}\NormalTok{(}\StringTok{"ORDER\_TS"}\NormalTok{).alias(}\StringTok{"LAST\_ORDER\_TS"}\NormalTok{),}
\NormalTok{    )}
\NormalTok{)}

\CommentTok{\# Create FeatureView with DataFrame{-}defined features}
\NormalTok{user\_purchase\_fv }\OperatorTok{=}\NormalTok{ FeatureView(}
\NormalTok{    name}\OperatorTok{=}\StringTok{"USER\_PURCHASE\_FEATURES"}\NormalTok{,}
\NormalTok{    entities}\OperatorTok{=}\NormalTok{[user\_entity],}
\NormalTok{    feature\_df}\OperatorTok{=}\NormalTok{user\_purchase\_df,  }\CommentTok{\# DataFrame defines the features}
\NormalTok{    timestamp\_col}\OperatorTok{=}\StringTok{"LAST\_ORDER\_TS"}\NormalTok{,}
\NormalTok{    refresh\_freq}\OperatorTok{=}\StringTok{"1 hour"}\NormalTok{,}
\NormalTok{)}
\end{Highlighting}
\end{Shaded}

\hypertarget{method-2-pure-sql-via-session.sql}{%
\subsubsection{Method 2: Pure SQL via
session.sql()}\label{method-2-pure-sql-via-session.sql}}

For complex queries or when you prefer SQL, use \texttt{session.sql()}
to define features:

\begin{quote}
📁 \textbf{Full code:}
\href{_code/feature_sql_api.py}{\texttt{\_code/feature\_sql\_api.py}}
\end{quote}

\begin{Shaded}
\begin{Highlighting}[]
\CommentTok{\# Define features using pure SQL}
\NormalTok{user\_features\_df }\OperatorTok{=}\NormalTok{ session.sql(}\StringTok{"""}
\StringTok{    SELECT }
\StringTok{        USER\_ID,}
\StringTok{        COUNT(DISTINCT ORDER\_ID) AS ORDER\_CNT,}
\StringTok{        SUM(AMOUNT) AS TOTAL\_SPEND,}
\StringTok{        AVG(AMOUNT) AS AVG\_ORDER\_AMT,}
\StringTok{        MAX(ORDER\_TS) AS LAST\_ORDER\_TS,}
\StringTok{        DATEDIFF(\textquotesingle{}day\textquotesingle{}, MIN(ORDER\_TS), MAX(ORDER\_TS)) AS CUSTOMER\_TENURE\_DAYS}
\StringTok{    FROM ORDERS}
\StringTok{    GROUP BY USER\_ID}
\StringTok{"""}\NormalTok{)}

\CommentTok{\# Create FeatureView from SQL{-}defined DataFrame}
\NormalTok{user\_purchase\_fv }\OperatorTok{=}\NormalTok{ FeatureView(}
\NormalTok{    name}\OperatorTok{=}\StringTok{"USER\_PURCHASE\_FEATURES"}\NormalTok{,}
\NormalTok{    entities}\OperatorTok{=}\NormalTok{[user\_entity],}
\NormalTok{    feature\_df}\OperatorTok{=}\NormalTok{user\_features\_df,  }\CommentTok{\# SQL query defines the features}
\NormalTok{    timestamp\_col}\OperatorTok{=}\StringTok{"LAST\_ORDER\_TS"}\NormalTok{,}
\NormalTok{    refresh\_freq}\OperatorTok{=}\StringTok{"1 hour"}\NormalTok{,}
\NormalTok{)}
\end{Highlighting}
\end{Shaded}

\begin{tcolorbox}[enhanced jigsaw, title=\textcolor{quarto-callout-tip-color}{\faLightbulb}\hspace{0.5em}{When to Use SQL}, coltitle=black, left=2mm, colback=white, toptitle=1mm, opacityback=0, colbacktitle=quarto-callout-tip-color!10!white, toprule=.15mm, colframe=quarto-callout-tip-color-frame, leftrule=.75mm, opacitybacktitle=0.6, breakable, bottomtitle=1mm, arc=.35mm, bottomrule=.15mm, titlerule=0mm, rightrule=.15mm]

The \texttt{session.sql()} approach is useful when migrating existing
SQL-based feature pipelines or when complex window functions are easier
to express in SQL.

\end{tcolorbox}

\hypertarget{sec-aggregations-api}{%
\subsubsection{Method 3: Feature Aggregation Class
(Time-Windowed)}\label{sec-aggregations-api}}

For time-windowed aggregations, the \texttt{Feature} class provides a
declarative API:

\begin{quote}
📁 \textbf{Full code:}
\href{_code/feature_aggregation_api.py}{\texttt{\_code/feature\_aggregation\_api.py}}
\end{quote}

\begin{Shaded}
\begin{Highlighting}[]
\ImportTok{from}\NormalTok{ snowflake.ml.feature\_store }\ImportTok{import}\NormalTok{ Feature}

\CommentTok{\# Define time{-}windowed aggregated features}
\NormalTok{features }\OperatorTok{=}\NormalTok{ [}
\NormalTok{    Feature.}\BuiltInTok{sum}\NormalTok{(}\StringTok{"AMOUNT"}\NormalTok{, }\StringTok{"7d"}\NormalTok{).alias(}\StringTok{"TOTAL\_SPEND\_7D"}\NormalTok{),}
\NormalTok{    Feature.count(}\StringTok{"ORDER\_ID"}\NormalTok{, }\StringTok{"30d"}\NormalTok{).alias(}\StringTok{"ORDER\_CNT\_30D"}\NormalTok{),}
\NormalTok{    Feature.avg(}\StringTok{"AMOUNT"}\NormalTok{, }\StringTok{"24h"}\NormalTok{).alias(}\StringTok{"AVG\_ORDER\_24H"}\NormalTok{),}
\NormalTok{    Feature.last\_n(}\StringTok{"PRODUCT\_ID"}\NormalTok{, }\StringTok{"7d"}\NormalTok{, n}\OperatorTok{=}\DecValTok{5}\NormalTok{).alias(}\StringTok{"RECENT\_PRODUCTS"}\NormalTok{),}
\NormalTok{]}

\CommentTok{\# Use in a tiled FeatureView for efficient computation}
\NormalTok{tiled\_fv }\OperatorTok{=}\NormalTok{ FeatureView(}
\NormalTok{    name}\OperatorTok{=}\StringTok{"USER\_PURCHASE\_AGGREGATES"}\NormalTok{,}
\NormalTok{    entities}\OperatorTok{=}\NormalTok{[user\_entity],}
\NormalTok{    feature\_df}\OperatorTok{=}\NormalTok{session.table(}\StringTok{"ORDERS"}\NormalTok{),}
\NormalTok{    timestamp\_col}\OperatorTok{=}\StringTok{"ORDER\_TS"}\NormalTok{,}
\NormalTok{    refresh\_freq}\OperatorTok{=}\StringTok{"1 hour"}\NormalTok{,}
\NormalTok{    feature\_granularity}\OperatorTok{=}\StringTok{"1 hour"}\NormalTok{,  }\CommentTok{\# Tile size for incremental computation}
\NormalTok{    features}\OperatorTok{=}\NormalTok{features,             }\CommentTok{\# Aggregated features}
\NormalTok{)}
\end{Highlighting}
\end{Shaded}

\begin{tcolorbox}[enhanced jigsaw, title=\textcolor{quarto-callout-note-color}{\faInfo}\hspace{0.5em}{🆕 New in snowflake-ml 1.21+}, coltitle=black, left=2mm, colback=white, toptitle=1mm, opacityback=0, colbacktitle=quarto-callout-note-color!10!white, toprule=.15mm, colframe=quarto-callout-note-color-frame, leftrule=.75mm, opacitybacktitle=0.6, breakable, bottomtitle=1mm, arc=.35mm, bottomrule=.15mm, titlerule=0mm, rightrule=.15mm]

The \texttt{Feature} aggregation class provides declarative syntax for
time-windowed aggregations with automatic tiling for efficient
incremental computation. See
\protect\hyperlink{chapter-07-aggregations-api}{Chapter 07: Aggregations
API} for the complete API reference.

\end{tcolorbox}

\hypertarget{choosing-the-right-method}{%
\subsubsection{Choosing the Right
Method}\label{choosing-the-right-method}}

\begin{longtable}[]{@{}
  >{\raggedright\arraybackslash}p{(\columnwidth - 4\tabcolsep) * \real{0.2667}}
  >{\raggedright\arraybackslash}p{(\columnwidth - 4\tabcolsep) * \real{0.3333}}
  >{\raggedright\arraybackslash}p{(\columnwidth - 4\tabcolsep) * \real{0.4000}}@{}}
\toprule\noalign{}
\begin{minipage}[b]{\linewidth}\raggedright
Method
\end{minipage} & \begin{minipage}[b]{\linewidth}\raggedright
Best For
\end{minipage} & \begin{minipage}[b]{\linewidth}\raggedright
Complexity
\end{minipage} \\
\midrule\noalign{}
\endhead
\bottomrule\noalign{}
\endlastfoot
\textbf{DataFrame API} & Simple aggregations, joins, filtering & Low \\
\textbf{SQL via session.sql()} & Complex queries, window functions,
existing SQL & Medium \\
\textbf{Feature class} & Time-windowed aggregations, sliding windows &
Advanced \\
\end{longtable}

\begin{center}\rule{0.5\linewidth}{0.5pt}\end{center}

\hypertarget{feature-slice}{%
\subsection{Feature Slice}\label{feature-slice}}

A \textbf{Feature Slice} identifies specific feature columns within a
FeatureView that you want to retrieve. This allows selective feature
retrieval rather than fetching all columns.

\begin{Shaded}
\begin{Highlighting}[]
\CommentTok{\# Get specific features from a FeatureView}
\NormalTok{user\_slice }\OperatorTok{=}\NormalTok{ user\_features\_fv.}\BuiltInTok{slice}\NormalTok{([}\StringTok{"TOTAL\_SPEND\_7D"}\NormalTok{, }\StringTok{"ORDER\_CNT\_30D"}\NormalTok{])}

\CommentTok{\# Use the slice in dataset generation}
\NormalTok{dataset }\OperatorTok{=}\NormalTok{ fs.generate\_dataset(}
\NormalTok{    spine\_df}\OperatorTok{=}\NormalTok{training\_spine,}
\NormalTok{    features}\OperatorTok{=}\NormalTok{[user\_slice],  }\CommentTok{\# Only retrieves sliced features}
\NormalTok{)}
\end{Highlighting}
\end{Shaded}

\begin{center}\rule{0.5\linewidth}{0.5pt}\end{center}

\hypertarget{sec-spine}{%
\section{The Spine: Connecting Features to ML}\label{sec-spine}}

The \textbf{Spine} is the foundation dataframe that defines which
entities need features and when. Its structure differs slightly between
training and inference:

\begin{longtable}[]{@{}lll@{}}
\toprule\noalign{}
Use Case & Required Columns & Optional Columns \\
\midrule\noalign{}
\endhead
\bottomrule\noalign{}
\endlastfoot
\textbf{Training} & Entity keys, Timestamp & Label (target variable) \\
\textbf{Batch Inference} & Entity keys, Timestamp & --- \\
\textbf{Online Inference} & Entity keys only & --- \\
\end{longtable}

\hypertarget{training-spine}{%
\subsubsection{Training Spine}\label{training-spine}}

For training, the spine defines the historical points where you want to
retrieve features, along with the target variable (label) you're
predicting:

\begin{quote}
📁 \textbf{Full code:}
\href{_code/spine_examples.py}{\texttt{\_code/spine\_examples.py}}
\end{quote}

\begin{Shaded}
\begin{Highlighting}[]
\CommentTok{\# Training spine: includes label for supervised learning}
\NormalTok{training\_spine }\OperatorTok{=}\NormalTok{ session.sql(}\StringTok{"""}
\StringTok{    SELECT }
\StringTok{        USER\_ID,                           {-}{-} Entity key}
\StringTok{        SUBSCRIPTION\_END\_TS AS EVENT\_TS,   {-}{-} Point{-}in{-}time timestamp  }
\StringTok{        CASE WHEN CHURNED THEN 1 ELSE 0 END AS LABEL  {-}{-} Target variable}
\StringTok{    FROM SUBSCRIPTION\_EVENTS}
\StringTok{    WHERE EVENT\_TYPE = \textquotesingle{}renewal\_decision\textquotesingle{}}
\StringTok{"""}\NormalTok{)}

\CommentTok{\# Generate training dataset with features joined to spine}
\NormalTok{training\_set }\OperatorTok{=}\NormalTok{ fs.generate\_dataset(}
\NormalTok{    spine\_df}\OperatorTok{=}\NormalTok{training\_spine,}
\NormalTok{    features}\OperatorTok{=}\NormalTok{[user\_features\_fv, session\_features\_fv],}
\NormalTok{    spine\_timestamp\_col}\OperatorTok{=}\StringTok{"EVENT\_TS"}\NormalTok{,}
\NormalTok{)}
\end{Highlighting}
\end{Shaded}

\hypertarget{batch-inference-spine}{%
\subsubsection{Batch Inference Spine}\label{batch-inference-spine}}

For batch inference, you only need entity keys and timestamps---there's
no label because that's what the model will predict:

\begin{Shaded}
\begin{Highlighting}[]
\CommentTok{\# Inference spine: no label column}
\NormalTok{inference\_spine }\OperatorTok{=}\NormalTok{ session.sql(}\StringTok{"""}
\StringTok{    SELECT }
\StringTok{        USER\_ID,                           {-}{-} Entity key}
\StringTok{        CURRENT\_TIMESTAMP() AS EVENT\_TS    {-}{-} Features as of now}
\StringTok{    FROM ACTIVE\_USERS}
\StringTok{"""}\NormalTok{)}

\CommentTok{\# Retrieve features for prediction}
\NormalTok{inference\_data }\OperatorTok{=}\NormalTok{ fs.generate\_dataset(}
\NormalTok{    spine\_df}\OperatorTok{=}\NormalTok{inference\_spine,}
\NormalTok{    features}\OperatorTok{=}\NormalTok{[user\_features\_fv, session\_features\_fv],}
\NormalTok{    spine\_timestamp\_col}\OperatorTok{=}\StringTok{"EVENT\_TS"}\NormalTok{,}
\NormalTok{)}
\end{Highlighting}
\end{Shaded}

\hypertarget{online-inference}{%
\subsubsection{Online Inference}\label{online-inference}}

For real-time serving, features are retrieved from Online Feature Tables
using entity keys only---no timestamp is needed since OFTs store only
the current (latest) feature values:

\begin{Shaded}
\begin{Highlighting}[]
\CommentTok{\# Online serving: entity keys only, retrieves current feature values}
\NormalTok{features }\OperatorTok{=}\NormalTok{ fs.retrieve\_feature\_values(}
\NormalTok{    spine\_df}\OperatorTok{=}\NormalTok{session.create\_dataframe([\{}\StringTok{"USER\_ID"}\NormalTok{: }\StringTok{"usr\_001"}\NormalTok{\}]),}
\NormalTok{    features}\OperatorTok{=}\NormalTok{[user\_features\_fv],}
\NormalTok{)}
\end{Highlighting}
\end{Shaded}

\hypertarget{how-spine-works}{%
\subsection{How Spine Works}\label{how-spine-works}}

\begin{verbatim}
┌──────────────────────────────────────────────────────────────────────────┐
│                      SPINE-BASED FEATURE RETRIEVAL                       │
├──────────────────────────────────────────────────────────────────────────┤
│                                                                          │
│  SPINE DataFrame                      FeatureViews                       │
│  ┌───────────────────────┐            ┌───────────────────────┐          │
│  │ USER_ID  │ EVENT_TS   │            │ USER_PURCHASE_FV      │          │
│  │──────────┼────────────│    JOIN    │  • TOTAL_SPEND_7D     │          │
│  │ usr_001  │ 2025-01-15 │───────────▶│  • ORDER_CNT_30D      │          │
│  │ usr_002  │ 2025-01-16 │            │  • AVG_ORDER_24H      │          │
│  │ usr_003  │ 2025-01-14 │            └───────────────────────┘          │
│  └───────────────────────┘                                               │
│            │                          ┌───────────────────────┐          │
│            │                          │ USER_SESSION_FV       │          │
│            └─────────────────────────▶│  • SESSION_CNT_7D     │          │
│                                       │  • AVG_SESSION_DUR    │          │
│                                       └───────────────────────┘          │
│                                                                          │
│  Result: Features computed as-of EVENT_TS for each USER_ID               │
└──────────────────────────────────────────────────────────────────────────┘
\end{verbatim}

The Feature Store joins each spine row to the appropriate FeatureViews,
ensuring features are computed using only data available at the
\texttt{EVENT\_TS} timestamp.

\begin{tcolorbox}[enhanced jigsaw, title=\textcolor{quarto-callout-note-color}{\faInfo}\hspace{0.5em}{📖 Deep Dive}, coltitle=black, left=2mm, colback=white, toptitle=1mm, opacityback=0, colbacktitle=quarto-callout-note-color!10!white, toprule=.15mm, colframe=quarto-callout-note-color-frame, leftrule=.75mm, opacitybacktitle=0.6, breakable, bottomtitle=1mm, arc=.35mm, bottomrule=.15mm, titlerule=0mm, rightrule=.15mm]

See \protect\hyperlink{chapter-10-training-inference}{Chapter 10:
Training \& Inference} for spine design patterns and best practices.

\end{tcolorbox}

\begin{center}\rule{0.5\linewidth}{0.5pt}\end{center}

\hypertarget{sec-architecture}{%
\section{Feature Store Architecture}\label{sec-architecture}}

Snowflake Feature Store leverages native Snowflake capabilities:

\begin{verbatim}
┌───────────────────────────────────────────────────────────────────────────┐
│                    SNOWFLAKE FEATURE STORE ARCHITECTURE                   │
├───────────────────────────────────────────────────────────────────────────┤
│                                                                           │
│  ┌────────────────┐      ┌────────────────────────────────────────────┐   │
│  │  Source Data   │      │           Feature Store Schema             │   │
│  │                │      │                                            │   │
│  │  • Tables      │      │  ┌──────────────────────────────────────┐  │   │
│  │  • Views       │─────▶│  │           FeatureViews               │  │   │
│  │  • Ext Tables  │      │  │  ┌───────────────┐ ┌──────────────┐  │  │   │
│  │  • Streams     │      │  │  │ Dynamic Table │ │    View      │  │  │   │
│  │                │      │  │  │ (materialized)│ │ (query-time) │  │  │   │
│  └────────────────┘      │  │  └───────────────┘ └──────────────┘  │  │   │
│                          │  └──────────────────────────────────────┘  │   │
│                          │                                            │   │
│                          │  ┌──────────────────────────────────────┐  │   │
│                          │  │         Metadata (Tags)              │  │   │
│                          │  │  • Entity definitions                │  │   │
│                          │  │  • Feature descriptions              │  │   │
│                          │  │  • Lineage information               │  │   │
│                          │  └──────────────────────────────────────┘  │   │
│                          └────────────────────────────────────────────┘   │
│                                            │                              │
│                          ┌─────────────────┼─────────────────┐            │
│                          │                 │                 │            │
│                          ▼                 ▼                 ▼            │
│                   ┌────────────┐   ┌────────────┐   ┌────────────┐        │
│                   │  Training  │   │   Batch    │   │   Online   │        │
│                   │  Datasets  │   │ Inference  │   │  Serving   │        │
│                   │            │   │            │   │   (OFT)    │        │
│                   └────────────┘   └────────────┘   └────────────┘        │
│                                                                           │
└───────────────────────────────────────────────────────────────────────────┘
\end{verbatim}

\hypertarget{physical-implementation}{%
\subsection{Physical Implementation}\label{physical-implementation}}

\begin{longtable}[]{@{}
  >{\raggedright\arraybackslash}p{(\columnwidth - 4\tabcolsep) * \real{0.3864}}
  >{\raggedright\arraybackslash}p{(\columnwidth - 4\tabcolsep) * \real{0.4091}}
  >{\raggedright\arraybackslash}p{(\columnwidth - 4\tabcolsep) * \real{0.2045}}@{}}
\toprule\noalign{}
\begin{minipage}[b]{\linewidth}\raggedright
Logical Concept
\end{minipage} & \begin{minipage}[b]{\linewidth}\raggedright
Snowflake Object
\end{minipage} & \begin{minipage}[b]{\linewidth}\raggedright
Purpose
\end{minipage} \\
\midrule\noalign{}
\endhead
\bottomrule\noalign{}
\endlastfoot
Feature Store & Schema + Tags & Container for all Feature Store
objects \\
Entity & Tag & Metadata defining join keys \\
FeatureView (materialized) & Dynamic Table & Pre-computed features with
full history \\
FeatureView (query-time) & View & On-demand computed features \\
Online Feature Table (OFT) & Hybrid Table & Low-latency serving,
\textbf{current values only} \\
Dataset & Table or View & Materialized training/inference data \\
\end{longtable}

\begin{tcolorbox}[enhanced jigsaw, title=\textcolor{quarto-callout-important-color}{\faExclamation}\hspace{0.5em}{Online Feature Tables Store Current Values Only}, coltitle=black, left=2mm, colback=white, toptitle=1mm, opacityback=0, colbacktitle=quarto-callout-important-color!10!white, toprule=.15mm, colframe=quarto-callout-important-color-frame, leftrule=.75mm, opacitybacktitle=0.6, breakable, bottomtitle=1mm, arc=.35mm, bottomrule=.15mm, titlerule=0mm, rightrule=.15mm]

Online Feature Tables store only the \textbf{latest (current) value} for
each feature per entity---they do not retain historical feature values.
For point-in-time historical retrieval (training, batch inference), use
the standard FeatureView backed by Dynamic Table or View.

\end{tcolorbox}

\begin{center}\rule{0.5\linewidth}{0.5pt}\end{center}

\hypertarget{sec-taxonomy}{%
\section{Transformation Taxonomy}\label{sec-taxonomy}}

Understanding where transformations should occur is critical for
maintainable ML systems. We categorize transformations into three types:

\begin{longtable}[]{@{}
  >{\raggedright\arraybackslash}p{(\columnwidth - 6\tabcolsep) * \real{0.1579}}
  >{\raggedright\arraybackslash}p{(\columnwidth - 6\tabcolsep) * \real{0.2895}}
  >{\raggedright\arraybackslash}p{(\columnwidth - 6\tabcolsep) * \real{0.2632}}
  >{\raggedright\arraybackslash}p{(\columnwidth - 6\tabcolsep) * \real{0.2895}}@{}}
\toprule\noalign{}
\begin{minipage}[b]{\linewidth}\raggedright
Type
\end{minipage} & \begin{minipage}[b]{\linewidth}\raggedright
Full Name
\end{minipage} & \begin{minipage}[b]{\linewidth}\raggedright
Location
\end{minipage} & \begin{minipage}[b]{\linewidth}\raggedright
Stored In
\end{minipage} \\
\midrule\noalign{}
\endhead
\bottomrule\noalign{}
\endlastfoot
\textbf{MIT} & Model-Independent Transformations & Feature Pipeline &
FeatureView \\
\textbf{MDT} & Model-Dependent Transformations & Training + Inference
Pipeline & Model Registry \\
\textbf{ODT} & On-Demand Transformations & Inference Time & Not
stored \\
\end{longtable}

\hypertarget{quick-decision-guide}{%
\subsection{Quick Decision Guide}\label{quick-decision-guide}}

\begin{verbatim}
Is it reusable across multiple models?
├── YES → MIT (store in FeatureView)
└── NO  → Does it depend on training data statistics?
          ├── YES → MDT (store with model)
          └── NO  → Does it depend on request context?
                    ├── YES → ODT (compute at inference)
                    └── NO  → Probably MIT
\end{verbatim}

\begin{tcolorbox}[enhanced jigsaw, title=\textcolor{quarto-callout-note-color}{\faInfo}\hspace{0.5em}{📖 Deep Dive}, coltitle=black, left=2mm, colback=white, toptitle=1mm, opacityback=0, colbacktitle=quarto-callout-note-color!10!white, toprule=.15mm, colframe=quarto-callout-note-color-frame, leftrule=.75mm, opacitybacktitle=0.6, breakable, bottomtitle=1mm, arc=.35mm, bottomrule=.15mm, titlerule=0mm, rightrule=.15mm]

See \href{./transformation_taxonomy.md}{Transformation Taxonomy: MIT vs
MDT vs ODT} for detailed examples, anti-patterns, and implementation
guidance.

\end{tcolorbox}

\begin{center}\rule{0.5\linewidth}{0.5pt}\end{center}

\hypertarget{terminology-mapping}{%
\section{Terminology Mapping}\label{terminology-mapping}}

Snowflake Feature Store terminology aligns with industry standards.
Here's how terms map across platforms:

\begin{longtable}[]{@{}
  >{\raggedright\arraybackslash}p{(\columnwidth - 6\tabcolsep) * \real{0.2157}}
  >{\raggedright\arraybackslash}p{(\columnwidth - 6\tabcolsep) * \real{0.3725}}
  >{\raggedright\arraybackslash}p{(\columnwidth - 6\tabcolsep) * \real{0.1569}}
  >{\raggedright\arraybackslash}p{(\columnwidth - 6\tabcolsep) * \real{0.2549}}@{}}
\toprule\noalign{}
\begin{minipage}[b]{\linewidth}\raggedright
Snowflake
\end{minipage} & \begin{minipage}[b]{\linewidth}\raggedright
Dowling/Hopsworks
\end{minipage} & \begin{minipage}[b]{\linewidth}\raggedright
Tecton
\end{minipage} & \begin{minipage}[b]{\linewidth}\raggedright
Description
\end{minipage} \\
\midrule\noalign{}
\endhead
\bottomrule\noalign{}
\endlastfoot
\textbf{Feature Store} & Feature Store & Feature Store & Central feature
repository \\
\textbf{Entity} & Entity & Entity & Business object with join keys \\
\textbf{FeatureView} & Feature Group & Feature View & Collection of
related features \\
\textbf{Dynamic Table FV} & Offline Store & Batch Feature View &
Pre-computed, stored features \\
\textbf{View FV} & - & On-Demand Feature View & Query-time computed
features \\
\textbf{Online Feature Table} & Online Store & Online Store &
Low-latency serving \\
\textbf{Spine} & Spine DataFrame & Spine & Request keys + timestamps \\
\textbf{Feature Slice} & Feature & Feature & Individual feature
column \\
\texttt{timestamp\_col} & Event Time & Timestamp Column & Point-in-time
reference \\
\texttt{refresh\_freq} & Materialization Schedule & TTL & Update
frequency \\
\texttt{generate\_dataset()} & Get Training Data & Get Dataset &
Training data generation \\
\texttt{retrieve\_feature\_values()} & Get Feature Values & Get Online
Features & Feature serving \\
\end{longtable}

\hypertarget{key-terminology-notes}{%
\subsection{Key Terminology Notes}\label{key-terminology-notes}}

\begin{enumerate}
\def\labelenumi{\arabic{enumi}.}
\item
  \textbf{FeatureView vs Feature Group}: Snowflake uses ``FeatureView''
  while Hopsworks uses ``Feature Group.'' Both represent a collection of
  features computed together.
\item
  \textbf{Dynamic Table vs Offline Store}: Both refer to
  pre-materialized features for batch training and inference.
\item
  \textbf{Online Feature Table}: Snowflake's implementation uses Hybrid
  Tables for low-latency serving, equivalent to an ``Online Store'' in
  other platforms.
\end{enumerate}

\begin{center}\rule{0.5\linewidth}{0.5pt}\end{center}

\hypertarget{summary}{%
\section{Summary}\label{summary}}

\begin{longtable}[]{@{}
  >{\raggedright\arraybackslash}p{(\columnwidth - 4\tabcolsep) * \real{0.2432}}
  >{\raggedright\arraybackslash}p{(\columnwidth - 4\tabcolsep) * \real{0.3243}}
  >{\raggedright\arraybackslash}p{(\columnwidth - 4\tabcolsep) * \real{0.4324}}@{}}
\toprule\noalign{}
\begin{minipage}[b]{\linewidth}\raggedright
Concept
\end{minipage} & \begin{minipage}[b]{\linewidth}\raggedright
Definition
\end{minipage} & \begin{minipage}[b]{\linewidth}\raggedright
Key Attributes
\end{minipage} \\
\midrule\noalign{}
\endhead
\bottomrule\noalign{}
\endlastfoot
\textbf{Feature Store} & Centralized feature management layer & Schema +
Tags in Snowflake \\
\textbf{Entity} & Business object with join keys & \texttt{name},
\texttt{join\_keys}, \texttt{desc} \\
\textbf{FeatureView} & Feature collection with transformations &
\texttt{entities}, \texttt{feature\_df}, \texttt{refresh\_freq} \\
\textbf{Feature} & Individual computed value & Defined in DataFrame or
\texttt{Feature} class \\
\textbf{Spine} & Request DataFrame & Entity keys + timestamp + labels \\
\textbf{MIT/MDT/ODT} & Transformation taxonomy & Where transformations
belong \\
\end{longtable}

\begin{center}\rule{0.5\linewidth}{0.5pt}\end{center}

\hypertarget{next-steps-1}{%
\section{Next Steps}\label{next-steps-1}}

Continue to \protect\hyperlink{chapter-02-design-organization}{Chapter
02: Design \& Organization} to learn how to structure your Feature Store
for scale and maintainability.

\part{Design \& Architecture}

\hypertarget{chapter-02-design-organization}{%
\chapter{Chapter 02: Design
Organization}\label{chapter-02-design-organization}}

Coming soon

\hfill\break

\hypertarget{overview-2}{%
\section{Overview}\label{overview-2}}

This chapter is under development. Content will be migrated from the
existing guide.

\begin{tcolorbox}[enhanced jigsaw, title=\textcolor{quarto-callout-warning-color}{\faExclamationTriangle}\hspace{0.5em}{Work in Progress}, coltitle=black, left=2mm, colback=white, toptitle=1mm, opacityback=0, colbacktitle=quarto-callout-warning-color!10!white, toprule=.15mm, colframe=quarto-callout-warning-color-frame, leftrule=.75mm, opacitybacktitle=0.6, breakable, bottomtitle=1mm, arc=.35mm, bottomrule=.15mm, titlerule=0mm, rightrule=.15mm]

This chapter is being converted to Quarto format. Check back soon for
updated content.

\end{tcolorbox}

\hypertarget{placeholder}{%
\section{Placeholder}\label{placeholder}}

See the original \url{index.md} for current content.

\hypertarget{chapter-03-entities-hierarchies}{%
\chapter{Chapter 03: Entities
Hierarchies}\label{chapter-03-entities-hierarchies}}

Coming soon

\hfill\break

\hypertarget{overview-3}{%
\section{Overview}\label{overview-3}}

This chapter is under development. Content will be migrated from the
existing guide.

\begin{tcolorbox}[enhanced jigsaw, title=\textcolor{quarto-callout-warning-color}{\faExclamationTriangle}\hspace{0.5em}{Work in Progress}, coltitle=black, left=2mm, colback=white, toptitle=1mm, opacityback=0, colbacktitle=quarto-callout-warning-color!10!white, toprule=.15mm, colframe=quarto-callout-warning-color-frame, leftrule=.75mm, opacitybacktitle=0.6, breakable, bottomtitle=1mm, arc=.35mm, bottomrule=.15mm, titlerule=0mm, rightrule=.15mm]

This chapter is being converted to Quarto format. Check back soon for
updated content.

\end{tcolorbox}

\hypertarget{placeholder-1}{%
\section{Placeholder}\label{placeholder-1}}

See the original \url{index.md} for current content.

\hypertarget{chapter-04-feature-views}{%
\chapter{Chapter 04: Feature Views}\label{chapter-04-feature-views}}

Coming soon

\hfill\break

\hypertarget{overview-4}{%
\section{Overview}\label{overview-4}}

This chapter is under development. Content will be migrated from the
existing guide.

\begin{tcolorbox}[enhanced jigsaw, title=\textcolor{quarto-callout-warning-color}{\faExclamationTriangle}\hspace{0.5em}{Work in Progress}, coltitle=black, left=2mm, colback=white, toptitle=1mm, opacityback=0, colbacktitle=quarto-callout-warning-color!10!white, toprule=.15mm, colframe=quarto-callout-warning-color-frame, leftrule=.75mm, opacitybacktitle=0.6, breakable, bottomtitle=1mm, arc=.35mm, bottomrule=.15mm, titlerule=0mm, rightrule=.15mm]

This chapter is being converted to Quarto format. Check back soon for
updated content.

\end{tcolorbox}

\hypertarget{placeholder-2}{%
\section{Placeholder}\label{placeholder-2}}

See the original \url{index.md} for current content.

\part{Feature Engineering}

\hypertarget{chapter-05-feature-pipelines}{%
\chapter{Chapter 05: Feature
Pipelines}\label{chapter-05-feature-pipelines}}

Coming soon

\hfill\break

\hypertarget{overview-5}{%
\section{Overview}\label{overview-5}}

This chapter is under development. Content will be migrated from the
existing guide.

\begin{tcolorbox}[enhanced jigsaw, title=\textcolor{quarto-callout-warning-color}{\faExclamationTriangle}\hspace{0.5em}{Work in Progress}, coltitle=black, left=2mm, colback=white, toptitle=1mm, opacityback=0, colbacktitle=quarto-callout-warning-color!10!white, toprule=.15mm, colframe=quarto-callout-warning-color-frame, leftrule=.75mm, opacitybacktitle=0.6, breakable, bottomtitle=1mm, arc=.35mm, bottomrule=.15mm, titlerule=0mm, rightrule=.15mm]

This chapter is being converted to Quarto format. Check back soon for
updated content.

\end{tcolorbox}

\hypertarget{placeholder-3}{%
\section{Placeholder}\label{placeholder-3}}

See the original \url{index.md} for current content.

\hypertarget{chapter-06-temporal-features}{%
\chapter{Chapter 06: Temporal
Features}\label{chapter-06-temporal-features}}

Coming soon

\hfill\break

\hypertarget{overview-6}{%
\section{Overview}\label{overview-6}}

This chapter is under development. Content will be migrated from the
existing guide.

\begin{tcolorbox}[enhanced jigsaw, title=\textcolor{quarto-callout-warning-color}{\faExclamationTriangle}\hspace{0.5em}{Work in Progress}, coltitle=black, left=2mm, colback=white, toptitle=1mm, opacityback=0, colbacktitle=quarto-callout-warning-color!10!white, toprule=.15mm, colframe=quarto-callout-warning-color-frame, leftrule=.75mm, opacitybacktitle=0.6, breakable, bottomtitle=1mm, arc=.35mm, bottomrule=.15mm, titlerule=0mm, rightrule=.15mm]

This chapter is being converted to Quarto format. Check back soon for
updated content.

\end{tcolorbox}

\hypertarget{placeholder-4}{%
\section{Placeholder}\label{placeholder-4}}

See the original \url{index.md} for current content.

\hypertarget{chapter-07-aggregations-api}{%
\chapter{Chapter 07: Aggregations
Api}\label{chapter-07-aggregations-api}}

Coming soon

\hfill\break

\hypertarget{overview-7}{%
\section{Overview}\label{overview-7}}

This chapter is under development. Content will be migrated from the
existing guide.

\begin{tcolorbox}[enhanced jigsaw, title=\textcolor{quarto-callout-warning-color}{\faExclamationTriangle}\hspace{0.5em}{Work in Progress}, coltitle=black, left=2mm, colback=white, toptitle=1mm, opacityback=0, colbacktitle=quarto-callout-warning-color!10!white, toprule=.15mm, colframe=quarto-callout-warning-color-frame, leftrule=.75mm, opacitybacktitle=0.6, breakable, bottomtitle=1mm, arc=.35mm, bottomrule=.15mm, titlerule=0mm, rightrule=.15mm]

This chapter is being converted to Quarto format. Check back soon for
updated content.

\end{tcolorbox}

\hypertarget{placeholder-5}{%
\section{Placeholder}\label{placeholder-5}}

See the original \url{index.md} for current content.

\part{Production}

\hypertarget{chapter-08-online-features}{%
\chapter{Chapter 08: Online Features}\label{chapter-08-online-features}}

Coming soon

\hfill\break

\hypertarget{overview-8}{%
\section{Overview}\label{overview-8}}

This chapter is under development. Content will be migrated from the
existing guide.

\begin{tcolorbox}[enhanced jigsaw, title=\textcolor{quarto-callout-warning-color}{\faExclamationTriangle}\hspace{0.5em}{Work in Progress}, coltitle=black, left=2mm, colback=white, toptitle=1mm, opacityback=0, colbacktitle=quarto-callout-warning-color!10!white, toprule=.15mm, colframe=quarto-callout-warning-color-frame, leftrule=.75mm, opacitybacktitle=0.6, breakable, bottomtitle=1mm, arc=.35mm, bottomrule=.15mm, titlerule=0mm, rightrule=.15mm]

This chapter is being converted to Quarto format. Check back soon for
updated content.

\end{tcolorbox}

\hypertarget{placeholder-6}{%
\section{Placeholder}\label{placeholder-6}}

See the original \url{index.md} for current content.

\hypertarget{chapter-09-preprocessing}{%
\chapter{Chapter 09: Preprocessing}\label{chapter-09-preprocessing}}

Coming soon

\hfill\break

\hypertarget{overview-9}{%
\section{Overview}\label{overview-9}}

This chapter is under development. Content will be migrated from the
existing guide.

\begin{tcolorbox}[enhanced jigsaw, title=\textcolor{quarto-callout-warning-color}{\faExclamationTriangle}\hspace{0.5em}{Work in Progress}, coltitle=black, left=2mm, colback=white, toptitle=1mm, opacityback=0, colbacktitle=quarto-callout-warning-color!10!white, toprule=.15mm, colframe=quarto-callout-warning-color-frame, leftrule=.75mm, opacitybacktitle=0.6, breakable, bottomtitle=1mm, arc=.35mm, bottomrule=.15mm, titlerule=0mm, rightrule=.15mm]

This chapter is being converted to Quarto format. Check back soon for
updated content.

\end{tcolorbox}

\hypertarget{placeholder-7}{%
\section{Placeholder}\label{placeholder-7}}

See the original \url{index.md} for current content.

\hypertarget{chapter-10-training-inference}{%
\chapter{Chapter 10: Training
Inference}\label{chapter-10-training-inference}}

Coming soon

\hfill\break

\hypertarget{overview-10}{%
\section{Overview}\label{overview-10}}

This chapter is under development. Content will be migrated from the
existing guide.

\begin{tcolorbox}[enhanced jigsaw, title=\textcolor{quarto-callout-warning-color}{\faExclamationTriangle}\hspace{0.5em}{Work in Progress}, coltitle=black, left=2mm, colback=white, toptitle=1mm, opacityback=0, colbacktitle=quarto-callout-warning-color!10!white, toprule=.15mm, colframe=quarto-callout-warning-color-frame, leftrule=.75mm, opacitybacktitle=0.6, breakable, bottomtitle=1mm, arc=.35mm, bottomrule=.15mm, titlerule=0mm, rightrule=.15mm]

This chapter is being converted to Quarto format. Check back soon for
updated content.

\end{tcolorbox}

\hypertarget{placeholder-8}{%
\section{Placeholder}\label{placeholder-8}}

See the original \url{index.md} for current content.

\part{Operations \& Advanced}

\hypertarget{chapter-11-operations}{%
\chapter{Chapter 11: Operations}\label{chapter-11-operations}}

Coming soon

\hfill\break

\hypertarget{overview-11}{%
\section{Overview}\label{overview-11}}

This chapter is under development. Content will be migrated from the
existing guide.

\begin{tcolorbox}[enhanced jigsaw, title=\textcolor{quarto-callout-warning-color}{\faExclamationTriangle}\hspace{0.5em}{Work in Progress}, coltitle=black, left=2mm, colback=white, toptitle=1mm, opacityback=0, colbacktitle=quarto-callout-warning-color!10!white, toprule=.15mm, colframe=quarto-callout-warning-color-frame, leftrule=.75mm, opacitybacktitle=0.6, breakable, bottomtitle=1mm, arc=.35mm, bottomrule=.15mm, titlerule=0mm, rightrule=.15mm]

This chapter is being converted to Quarto format. Check back soon for
updated content.

\end{tcolorbox}

\hypertarget{placeholder-9}{%
\section{Placeholder}\label{placeholder-9}}

See the original \url{index.md} for current content.

\hypertarget{chapter-12-advanced-patterns}{%
\chapter{Chapter 12: Advanced
Patterns}\label{chapter-12-advanced-patterns}}

Coming soon

\hfill\break

\hypertarget{overview-12}{%
\section{Overview}\label{overview-12}}

This chapter is under development. Content will be migrated from the
existing guide.

\begin{tcolorbox}[enhanced jigsaw, title=\textcolor{quarto-callout-warning-color}{\faExclamationTriangle}\hspace{0.5em}{Work in Progress}, coltitle=black, left=2mm, colback=white, toptitle=1mm, opacityback=0, colbacktitle=quarto-callout-warning-color!10!white, toprule=.15mm, colframe=quarto-callout-warning-color-frame, leftrule=.75mm, opacitybacktitle=0.6, breakable, bottomtitle=1mm, arc=.35mm, bottomrule=.15mm, titlerule=0mm, rightrule=.15mm]

This chapter is being converted to Quarto format. Check back soon for
updated content.

\end{tcolorbox}

\hypertarget{placeholder-10}{%
\section{Placeholder}\label{placeholder-10}}

See the original \url{index.md} for current content.

\hypertarget{chapter-13-migration-guide}{%
\chapter{Chapter 13: Migration Guide}\label{chapter-13-migration-guide}}

Coming soon

\hfill\break

\hypertarget{overview-13}{%
\section{Overview}\label{overview-13}}

This chapter is under development. Content will be migrated from the
existing guide.

\begin{tcolorbox}[enhanced jigsaw, title=\textcolor{quarto-callout-warning-color}{\faExclamationTriangle}\hspace{0.5em}{Work in Progress}, coltitle=black, left=2mm, colback=white, toptitle=1mm, opacityback=0, colbacktitle=quarto-callout-warning-color!10!white, toprule=.15mm, colframe=quarto-callout-warning-color-frame, leftrule=.75mm, opacitybacktitle=0.6, breakable, bottomtitle=1mm, arc=.35mm, bottomrule=.15mm, titlerule=0mm, rightrule=.15mm]

This chapter is being converted to Quarto format. Check back soon for
updated content.

\end{tcolorbox}

\hypertarget{placeholder-11}{%
\section{Placeholder}\label{placeholder-11}}

See the original \url{index.md} for current content.

\cleardoublepage
\phantomsection
\addcontentsline{toc}{part}{Appendices}
\appendix

\hypertarget{appendix-a-sample-data}{%
\chapter{Appendix A: Sample Data}\label{appendix-a-sample-data}}

Clickstream dataset and public ML datasets

\hfill\break

\hypertarget{overview-14}{%
\section{Overview}\label{overview-14}}

This appendix documents the sample data used throughout the guide.

\begin{tcolorbox}[enhanced jigsaw, title=\textcolor{quarto-callout-warning-color}{\faExclamationTriangle}\hspace{0.5em}{Work in Progress}, coltitle=black, left=2mm, colback=white, toptitle=1mm, opacityback=0, colbacktitle=quarto-callout-warning-color!10!white, toprule=.15mm, colframe=quarto-callout-warning-color-frame, leftrule=.75mm, opacitybacktitle=0.6, breakable, bottomtitle=1mm, arc=.35mm, bottomrule=.15mm, titlerule=0mm, rightrule=.15mm]

This appendix is being converted to Quarto format.

\end{tcolorbox}

See the original \url{index.md} for current content.



\end{document}
